% Combined drivers resolve supplement labels; standalone main uses descriptive references.
\providecommand{\KBSuppRef}[3]{#2~\ref{#1}}
% Shared scientific narrative. Audit material is in supplement.tex.

% Shared scientific main text for the maintained compact and anonymous TMLR layouts.
% Historical numerical authorities are retained. Completed next-phase comparisons
% have separate hash-bound authorities; supplementary material follows the bibliography.

\section{Introduction}
\label{sec:compact-intro}

An update that improves a model on one deployment distribution can degrade it on another.
Test-time adaptation proposes updates using unlabeled data through entropy minimization,
filtering, regularization, and related methods~\cite{wang2021tent,niu2022eata,niu2023sar,wang2022cotta}.
Choosing whether to serve the resulting candidate is a separate decision. The challenge is to
justify that choice when the target labels needed to measure its benefit are unavailable.

Our primary research question is:
\emph{When can evidence available without new target labels justify replacing a frozen predictor
with a fixed adapted candidate?}
Let $f_0$ be the frozen predictor and $f_a$ the candidate. Their population benefit is
\begin{equation}
  \Delta=R_T(f_0)-R_T(f_a).
  \label{eq:compact-benefit}
\end{equation}
A positive benefit favors the update; a negative benefit favors the frozen model. We distinguish
three conclusions: \adapt supports positive benefit, \freeze supports negative benefit, and
\abstain records insufficient evidence for either strict direction. Both non-adapt actions retain
$f_0$, but only freeze asserts that the candidate is worse.

K-Bound studies this decision as a partial-identification problem. In a binary model, we hold the
input law, predictors, and label-free evidence fixed and vary the target labels within a declared
class. The resulting feasible benefit set gives an exact boundary for a uniformly justified strict
decision. The bound on the hidden calibration residual is an assumption about that class, not a
quantity that an unlabeled batch supplies for free. A second result makes this information limit
explicit: an audit using the same evidence must still cover the unresolved residual range.

The theory asks what a fixed label-free evidence law and a declared residual bound can identify.
Knowability-Guided Adaptation (\KGA) studies the same update decision using additional information:
labeled historical shift outcomes. It fits an empirical benefit predictor and calibrates a
three-way decision rule. The connection is the decision problem; the empirical gate does not
implement the population frontier. Section~\ref{sec:compact-bridge} gives the full relationship
between their targets and assumptions.

\subsection{Three Scientific Contributions}

\begin{enumerate}
  \item \textbf{An exact benefit frontier on a declared class.}
  We construct target label laws with identical label-free evidence, derive their feasible benefit
  interval, and characterize its strict positive, strict negative, and unresolved regions,
  including the zero-benefit boundary.
  \item \textbf{An evidence-fibre audit limit.}
  We identify the least residual bound that the same observable law can support uniformly,
  explaining why historical labels help only through a stated link to the deployment domain.
  \item \textbf{An empirical paired-benefit gate.}
  We give a reproducible three-way interval procedure and evaluate its commitment behavior under
  controlled and natural shifts. The standard coverage-to-action implication states exactly what
  follows if the interval covers its declared target.
\end{enumerate}

The empirical studies then ask whether an interval improves the update decision and whether
its calibration transfers beyond the fitting conditions. Controlled CIFAR-10-C results are
candidate-dependent: Tent and EATA have lower point regret than either fixed policy, whereas
SAR favors always-adapt. A comparison using a shared point predictor isolates the contribution
of calibration from the predictor itself (Section~\ref{sec:nextphase-calibration}). It finds a
benefit for CIFAR Tent at one harm-weighted operating point, but simple rules win elsewhere
and nominal interval inclusion deteriorates on held-out corruptions.

Natural-shift studies expose a further distinction between retaining utility and selecting
useful updates. On CCT-20, KGA retained the frozen model in all 45 cells, avoiding the measured
loss of always-adapt while making no ADAPT decisions. The So2Sat studies stopped before target
scoring. These findings leave useful prospective routing unresolved. A separate paired-transport
construction shows how a declared conditional-shift assumption can support a finite-sample
multiclass decision, with synthetic examples illustrating its value and its failure when that
assumption is violated. Additional adverse and incomplete studies are reported in the supplement.

\begin{figure*}[t]
  \centering
  \kbgraphics[width=0.94\textwidth]{figures/fig_decision_flow.pdf}
  \caption{KGA's empirical workflow. A shadow candidate proposes an update; a learned benefit
  prediction and residual interval determine whether to accept it. Historical labels support
  development and calibration, while the new decision uses no evaluation labels.}
  \label{fig:compact-flow}
\end{figure*}

\section{Related Work}
\label{sec:compact-related}

\subsection{Candidate Adaptation and Update Selection}

Tent, EATA, SAR, CoTTA, MEMO, SHOT, and test-time training improve candidate construction through
entropy objectives, sample filtering, regularization, temporal stabilization, augmentation, or
auxiliary tasks~\cite{wang2021tent,niu2022eata,niu2023sar,wang2022cotta,zhang2022memo,
liang2020shot,sun2020ttt}. Selective methods also choose where or when to update. TTSA selects
modality-specific adapters~\cite{chen2025ttsa}; GALA selects layers and filters samples using
gradient alignment~\cite{sahoo2025gala}; Informed Adaptation learns topological features associated
with useful updates~\cite{mutlu2026informed}. KGA evaluates a fixed proposed candidate against its
frozen baseline. Learned routing and rollback are established ideas; their existence is not our
novelty claim.
MORPHEUS~\cite{danilowski2026morpheus}, presented at the ICLR 2026 Test-Time Updates
workshop, learns historical regressors from source-model entropy and embedding
geometry to predict post-adaptation accuracy and choose a method before executing
it. This directly overlaps the practical motivation of avoiding harmful updates.
KGA's proposed distinction is the calibrated paired-benefit operating point for
one executed shadow candidate; whether that calibration improves useful decisions
over a shared point predictor is an empirical question evaluated separately.

\subsection{Label-Free Performance and Strategy Estimation}

ATC, agreement-on-the-line, AETTA, and prediction-matrix statistics estimate performance from
unlabeled outputs and suitable development information
~\cite{garg2022atc,baek2022aol,miller2021aol,lee2024aetta,deng2023confidence,xie2024mano}.
The closest empirical comparison is Kim et al.'s NeurIPS 2024 work,
\emph{Test-Time Adaptation Induces Stronger Accuracy and Agreement-on-the-Line}
~\cite{kim2024ttaline}. It uses agreement-based ALine estimators after adaptation, supports
hyperparameter selection, and compares adaptation strategies without labeled OOD data.

KGA instead fits the paired observed benefit of one frozen/candidate pair and uses an interval
to make a three-way decision. This is a difference in the fitted target and decision procedure,
not a claim that accuracy estimates cannot be differenced to compare models. In Section~\ref{sec:synchronized-baselines}
(Table~\ref{tab:synchronized-baselines}), we compare current KGA with fixed
policies, scalar evidence rules and a point gate sharing its exact predictions.
These saved-feature comparisons isolate decision rules; they do not establish equivalence to
complete published TTA/ALine, AETTA, or POEM systems. Method-specific implementation and
calibration qualifications are given in \KBSuppRef{app:baseline-audit}{Appendix}{the supplementary baseline audit}.

\subsection{Monitoring and Adaptation Controllers}

POEM protects adaptation through entropy matching and online betting~\cite{bar2024poem}.
Schirmer et al. monitor the running risk of an adapting model relative to a source-risk reference,
under a proxy-informativeness condition~\cite{schirmer2025monitoring}; the supervised precursor
uses a labeled test stream~\cite{podkopaev2022tracking}. Drift-to-Action Controllers use delayed
labels and online certificates to gate adaptation, abstention, rollback, and retraining
~\cite{lamaakal2026drifttoaction}. These methods differ in the risk target, label access, and
sequential guarantee. KGA concerns one paired benefit and does not establish an anytime-valid
risk guarantee.

\subsection{Partial Identification and Robust Decisions}

An identified set describes the values compatible with the observations and assumptions.
Decision theory under partial identification studies how to act when that set leaves payoffs
ambiguous; Christensen et al. combine worst-case payoff uncertainty with statistical learning of
identified parameters~\cite{christensen2026optimal}. K-Bound addresses the narrower question of
a uniformly correct strict benefit sign. It does not solve their general optimal-decision problem
or claim that interval thresholding itself is new.

Indistinguishable adaptation problems are already central to domain-adaptation impossibility
~\cite{bendavid2010theory,bendavid2010impossibility}. K-Bound's specific contribution is a
disagreement-conditional label-kernel construction, its clipped benefit set, and the associated
residual-audit floor. The proof does not introduce label-free impossibility.
Jhawar and Wang~\cite{jhawar2026identifiable} independently study when an unlabeled
observation channel identifies the best TTA action. Their September 2026 preprint
includes a finite-batch Gaussian keep/recenter boundary and benchmark examples
where stream structure changes the preferred action. Thus action identification
being weaker than mechanism identification is not a unique claim of this paper.
Our population result fixes a realized predictor pair and a disagreement-conditional
residual class, deriving its clipped paired-benefit set and strict-sign boundary.
The distinction lies in that explicit compatible-law construction and its audit
requirements, not in the general fibre-consistency principle.
Risk-estimation assumptions based on conditional independence~\cite{steinhardt2016}, label shift
~\cite{lipton2018bbse,garg2020labelshift}, or disagreement discrepancy~\cite{rosenfeld2023dis2}
restrict the compatible laws. Likewise, MMD-credal risk intervals impose RKHS and covariate-shift
structure~\cite{ariq2026mmd}. Such restrictions can change the identified set; they are not
contradicted by a converse theorem for a richer class.
Li, Han, and Ma~\cite{li2026matrixconstraints} construct finite-sample label-shift
confidence regions through matrix constraints and linear programming. Consequently,
optimizing a compatible probability table is established machinery. The separate
paired-transport extension evaluated here targets the benefit contrast directly,
retains joint prediction-pair information, and reports sensitivity to an externally
specified conditional-transport violation. It does not estimate that violation
from unlabeled data or claim a new general-purpose LP confidence method.

\subsection{Marginal and Conditional Coverage}

Selective classification and reject-option learning abstain on individual predictions
~\cite{geifman2017selective,chow1970reject,cortes2016boosting}. KGA abstains on an update.
Its coverage-to-action proposition applies standard interval logic
~\cite{tibshirani2019conformal,barber2023beyond,bates2021rcps,angelopoulos2025ltt}.
The main statistical work is to justify coverage, especially after selection or under shift.

Jin and Ren's JOMI procedure targets coverage conditional on a unit being selected, under
exchangeability and a compatible selection rule~\cite{jin2025selection}. Gibbs et al. provide
conditional guarantees over specified classes of shifts using a different calibration construction
~\cite{gibbs2025conditional}. KGA's ordinary residual interval implements neither procedure.
Its marginal false-adapt bound therefore must be distinguished from error conditional on ADAPT
and from coverage within each target environment. Learn Then Test separately addresses
configuration selection and multiplicity~\cite{angelopoulos2025ltt}.

\begin{table*}[t]
\centering
\footnotesize
\caption{Closest roles, targets, and remaining distinctions. Literature comparisons are conceptual here;
a historical AETTA comparison is reported in \KBSuppRef{app:baseline-audit}{Appendix}{the supplementary baseline audit}, as a descriptive comparison whose rank-clipped calibration does not establish nominal coverage, while native TTA/ALine and full POEM comparisons remain incomplete.}
\label{tab:related-comparison}
\begin{tabular}{@{}p{0.20\textwidth}p{0.27\textwidth}p{0.43\textwidth}@{}}
\toprule
Approach & Main target & Relation to this paper \\
\midrule
TTA/ALine; AETTA & adapted-model performance and selection & closest empirical estimators; their scores can inform paired comparisons \\
POEM; risk monitors & protected adaptation or running risk & different monitored target and sequential assumptions \\
Partial-identification decisions & uncertain payoffs across compatible laws & context for the population benefit set and robust sign decision \\
Selection/shift-conditional conformal & coverage after selection or within a shift class & stronger, specialized targets not supplied by KGA's marginal interval \\
K-Bound / KGA & feasible population benefit / predicted cell benefit & shared update question, distinct uncertainty objects and guarantees \\
\bottomrule
\end{tabular}
\end{table*}

\section{Setup, Notation, and Claim Levels}
\label{sec:compact-setup}

\begin{table*}[t]
\centering
\footnotesize
\caption{Notation and probability targets. Population and measured-cell quantities are kept separate.}
\label{tab:compact-bridge}
\begin{tabular}{@{}p{0.19\textwidth}p{0.16\textwidth}p{0.55\textwidth}@{}}
\toprule
Quantity & Level & Meaning \\
\midrule
$f_0,f_a$ & both & frozen predictor and realized shadow candidate \\
$\Delta$ & population & expected risk reduction $R_T(f_0)-R_T(f_a)$ \\
$D,\ \mu_T(D)$ & population & disagreement region and its target probability \\
$M$ & population & label-free mean correctness-score margin on $D$ \\
$\gamma,\ \beta$ & population & hidden calibration residual and its externally declared absolute bound \\
$\Delta^{\mathrm{cell}}$ & empirical & observed score difference on one predeclared evaluation unit \\
$\widehat\Delta,\ \varepsilon$ & empirical & predicted cell benefit and residual-based interval radius \\
$\mathrm{FA}_{\mathrm u}$ & empirical & probability of ADAPT and a nonpositive cell benefit \\
$\mathrm{FA}_{\mathrm c}$ & empirical & probability of nonpositive cell benefit conditional on ADAPT \\
Commitment rate & empirical & fraction of ADAPT or FREEZE decisions; not interval coverage \\
\bottomrule
\end{tabular}
\end{table*}

\subsection{Model Pair, Evidence, and Evaluation Unit}

Let $P_S$ and $P_T$ be the source and target laws and
$R_T(f)=\mathbb E_{(X,Y)\sim P_T}[\ell(f(X),Y)]$.
Population risks integrate a fresh draw from $P_T$, holding the realized model pair fixed.
An adapter $\mathcal A$ proposes $f_a$ from $f_0$ and unlabeled inputs $X_{1:m}$.
The decision can read
\[
  Z=\phi(X_{1:m},f_0,f_a),
\]
including predictions, logits, and update statistics, but not the new evaluation labels.

For evaluation unit $j$, with records $\mathcal E_j$ and a fixed higher-is-better score $S$, define
\[
  \Delta_j^{\mathrm{cell}}=S(f_a;\mathcal E_j)-S(f_0;\mathcal E_j).
\]
The two models are scored on the same records. This measured quantity is not population risk.
For accuracy it is a sample average of paired correctness differences; for macro-F1 it is a
difference of dataset-level scores, not an average of per-example losses. Metric scales are
reported separately.

\subsection{Actions and Their Operational Meaning}

\begin{definition}[Decision semantics]
\adapt commits $f_a$ when the interval supports positive benefit.
\freeze rejects the candidate when it supports negative benefit.
\abstain retains $f_0$ when neither strict direction is supported or the assessment is unavailable.
\end{definition}

Freeze is a supported rejection. Abstention is an unresolved assessment. An operator can attach
monitoring, label acquisition, or human review to the abstention record; the experiments here
evaluate only the frozen fallback and uncertainty log, not those additional interventions.
Retaining $f_0$ after missing evidence must be logged as ABSTAIN, not as a certified FREEZE.

The observed-cell oracle chooses the larger score. Regret is the policy's score shortfall from
that oracle, averaged over the named units. Always-freeze regret is
$\max(\Delta_j^{\mathrm{cell}},0)$; always-adapt regret is
$\max(-\Delta_j^{\mathrm{cell}},0)$. KGA uses the candidate score only after ADAPT. Lower regret is
better, and a zero-benefit tie carries zero regret for either model even though it supports
neither strict sign.

\subsection{Commitment, Interval Coverage, and Error}

These are three different measurements:
\begin{align}
  \mathrm{Commit}
    &=\Pr(g\in\{\adapt,\freeze\}), \label{eq:commitment-rate}\\
  \mathrm{Cov}_{\mathrm{cell}}
    &=\Pr(|\widehat\Delta-\Delta^{\mathrm{cell}}|\le\varepsilon), \label{eq:interval-coverage}\\
  \mathrm{FA}_{\mathrm u}
    &=\Pr(g=\adapt,\Delta^{\mathrm{cell}}\le0). \label{eq:compact-fau}
\end{align}
Conditional false adaptation is
$\mathrm{FA}_{\mathrm c}=\Pr(\Delta^{\mathrm{cell}}\le0\mid g=\adapt)$, defined only when
$\Pr(g=\adapt)>0$. The operational proposition below bounds the marginal error under interval
coverage; it does not bound $\mathrm{FA}_{\mathrm c}$. A gate that never adapts has zero false
adaptations mechanically, so every error report includes action exposure.
Analogously, a false FREEZE is a FREEZE decision with $\Delta^{\mathrm{cell}}\ge0$; its
conditional frequency divides by FREEZE decisions and is undefined when that action is absent.

\begin{table*}[t]
\centering
\footnotesize
\caption{Three claim levels and their validity obligations. This table collects the distinctions
used throughout the paper.}
\label{tab:compact-validity}
\begin{tabular}{@{}p{0.18\textwidth}p{0.34\textwidth}p{0.39\textwidth}@{}}
\toprule
Claim level & Required premise & Conclusion and boundary \\
\midrule
Population frontier & binary $0/1$ loss, fixed augmented evidence, declared residual class; richness for necessity & exact strict-sign frontier within that class, not a learned residual bound \\
Operational proposition & marginal interval coverage for a named scalar target and unit & marginal false commitment for that target, not automatic conditional or anytime control \\
Empirical evaluation & fixed recorded candidates, splits, metric and action rule & measured cell regret, interval inclusion and action exposure; population transfer needs a separate argument \\
\bottomrule
\end{tabular}
\end{table*}

The population theorem uses binary labels and $0/1$ loss. The empirical gate can instead target
any predeclared scalar score difference, including multiclass metrics. This interface does not
extend the binary theorem. In general multiclass classification,
$\Delta=\Pr_T(D)(p_a-p_0)$, with $p_a=\Pr_T(f_a=Y\mid D)$ and
$p_0=\Pr_T(f_0=Y\mid D)$; only the binary setting forces $p_0=1-p_a$.

\section{Population Partial-Identification Frontier}
\label{sec:compact-theory}

The result describes which benefit signs are compatible with a fixed label-free evidence law
and a declared bound on the hidden residual.
The population analysis uses binary classification and $0/1$ loss, where disagreement implies that exactly one predictor is correct. This permits the exact benefit reduction and strict-commitment frontier derived below. The impossibility result applies within the declared label-kernel class; it does not rule out certification under additional structural restrictions or establish the same frontier for general multiclass problems. \KGA separately evaluates an empirical scalar-benefit gate, including on multiclass benchmarks.

\subsection{Disagreement-Region Model}

\begin{assumption}[Deployment setup]
\label{ass:deploy}
Let $Y\in\{0,1\}$, use $0/1$ loss, fix measurable binary predictors $f_0,f_a$, and assume
$\mu_T(D)>0$ for $D=\{x:f_0(x)\ne f_a(x)\}$. Let
$s:\mathcal X\to[0,1]$ be a fixed measurable candidate-correctness score, fitted using source
information, and let the target law admit the measurable pointwise correctness kernel
$\eta_a(x)=\Pr_T(f_a(X)=Y\mid X=x)$. On $D$, set
\begin{equation}
 M=\E[s(X)\mid D]-\tfrac12,
 \qquad
 \gamma=\E[\eta_a(X)-s(X)\mid D].
 \label{eq:compact-Mgamma}
\end{equation}
The declared population class is
$\mathcal C_\beta=\{P_T:|\gamma(P_T)|\le\beta\}$ for an externally supplied
$\beta\ge0$.
Because $s\in[0,1]$, the feasible margin satisfies $M\in[-\tfrac12,\tfrac12]$.
The identities and frontier below do not require a source-calibration property of $s$.
\end{assumption}

\begin{remark}[Residual identity]
\label{rem:gamma-residual}
The definition gives
\begin{align*}
M+\gamma
&=\bigl(\E[s(X)\mid D]-\tfrac12\bigr)
  +\E[\eta_a(X)-s(X)\mid D]\\
&=\E[\eta_a(X)\mid D]-\tfrac12.
\end{align*}
This is an identity, not empirical evidence that calibration transfers. We call $\gamma$ a
target disagreement-conditional calibration residual, not automatically source-to-target
drift. Ordinary source calibration need not hold after conditioning on $D$, so $\gamma$ can be
nonzero even when $P_T=P_S$. A literal label-kernel shift interpretation would require an
additional reference, for example the pointwise source correctness
$s(x)=\Pr_S(f_a(X)=Y\mid X=x)$. The no-shift counterexample is given in
\KBSuppRef{app:theorem-dependencies}{Appendix}{the supplementary theorem-dependency map}.
Calibration notions also depend on the conditioning information: top-label calibration, for
example, additionally conditions on the predicted class~\cite{gupta2021toplabel}.
Neither score nor top-label calibration automatically supplies calibration conditional on $D$.
\end{remark}

\begin{definition}[Risk alignment]
\label{def:risk-align}
Evidence $Z$ is risk-aligned on a class $\mathcal P$ if no two laws in $\mathcal P$ with opposite
nonzero benefit signs induce the same law of $Z$. This is weaker than uniform strict-sign
identification: a fibre containing zero and positive benefit is risk-aligned by this definition,
but does not support a uniformly positive claim. Population frontier claims take $M$ to be fixed on each
augmented evidence fibre, equivalently by using $\widetilde Z=(Z,M)$.
\end{definition}

\begin{definition}[Strict directional soundness]
\label{def:strict-sound}
At a fixed augmented-evidence law, \adapt is uniformly directionally sound when $\Delta(P)>0$ for
every compatible $P\in\mathcal C_\beta$; \freeze is sound when $\Delta(P)<0$ for every such law.
A pointwise maximal three-way rule commits whenever one strict action is sound and abstains
otherwise. The strict convention treats $\Delta=0$ as insufficient for either directional claim,
although the actions then have equal task risk.
\end{definition}

\subsection{Label-Kernel Freedom on the Disagreement Region}

\begin{lemma}[Label-kernel freedom on disagreement]
\label{lem:fibre}
Fix $(\mu_T,P_S,f_0,f_a,s)$, hence $D$ and $M$, and fix the conditional label law outside $D$.
For every measurable $\eta:D\to[0,1]$ there is a target law with input marginal $\mu_T$ and
$\Pr(f_a(X)=Y\mid X=x)=\eta(x)$ on $D$. Every label-free observable formed from the input batch,
the fixed models, and their outputs has the same law for all such $\eta$, while
\[
 \gamma=\E[\eta(X)\mid D]-\tfrac12-M.
\]
\end{lemma}
\begin{proof}
On $D$, binary predictors disagree, so their outputs are $0$ and $1$. Assign conditional mass
$\eta(x)$ to $Y=f_a(x)$ and $1-\eta(x)$ to $Y=f_0(x)$, retaining the fixed kernel outside $D$.
This changes no input marginal or model output and therefore changes no label-free evidence law.
The expression for $\gamma$ follows from Eq.~\eqref{eq:compact-Mgamma}.
\end{proof}

% Shared statements and proofs for the reduction, interior impossibility,
% closed-band boundary case, and exact strict-commitment frontier.
% Core load-bearing theory for the short-paper main text (with proofs).

\subsection{Benefit Reduction and Matched-Evidence Impossibility}

The first result connects the decision target to the three population symbols. Agreement outside
$D$ contributes no risk difference. Binary complementarity on $D$ then makes the sign of benefit
depend only on $M+\gamma$.

\begin{lemma}[Disagreement-region reduction]\label{lem:reduction}
Under Assumption~\ref{ass:deploy}, with $\bar a:=\Pr_T(f_a{=}Y\mid D)$,
\begin{align*}
\Delta &= 2\,\mu_T(D)\bigl(\bar a-\tfrac12\bigr),\\
\sign\Delta &= \sign(\bar a-\tfrac12)=\sign(M+\gamma).
\end{align*}
The identity $\sign\Delta=\sign(M+\gamma)$ holds for \emph{every} target law; it does not
assume $|\gamma|\le\beta$.
\end{lemma}
\begin{proof}
Outside $D$, the two predictors agree and incur the same $0/1$ loss. On $D$, exactly one is
correct. The adapted and frozen accuracies conditional on $D$ are therefore $\bar a$ and
$1-\bar a$. Hence $\Delta=2\mu_T(D)(\bar a-\tfrac12)$. The definitions of $M$ and $\gamma$ give
\begin{align*}
M+\gamma
&=\Bigl(\E_{\mu_T\mid D}[s]-\tfrac12\Bigr)
  +\E_{\mu_T\mid D}[\eta_a-s]\\
&=\E_{\mu_T\mid D}[\eta_a]-\tfrac12
 =\bar a-\tfrac12.
\end{align*}
Because $\mu_T(D)>0$, the three quantities have the same sign.
\end{proof}

The reduction tells us which hidden quantity controls the sign. The harder question is
identification. Can two target worlds agree on every permitted observation while assigning
opposite signs to $\Delta$? Inside the declared band, the answer is yes.

\begin{theorem}[Interior matched-evidence impossibility]\label{lem:nonid}
Fix $\beta>0$. For any feasible margin $M\in[-\tfrac12,\tfrac12]$ with $|M|<\beta$, there exist
$P_T^+,P_T^-\in\mathcal{C}_\beta$ with the following properties:
\begin{gather*}
  \Law(\widetilde Z\mid P_T^+)=\Law(\widetilde Z\mid P_T^-),\\
  M(P_T^+)=M(P_T^-)=M,\\
  \sign\Delta_+=-\sign\Delta_-\ne0,
\end{gather*}
where $\widetilde Z=(Z,M)$. If $\beta=0$, then $\mathcal C_0$ forces $\gamma=0$ and
$\sign\Delta=\sign M$; matched-evidence ambiguity therefore requires $\beta>0$.
\end{theorem}
\begin{proof}
Choose $0<\delta<\min\{\beta-|M|,\tfrac12\}$. For
$\theta\in\{+1,-1\}$, apply Lemma~\ref{lem:fibre} with the constant correctness kernel
\[
  \eta_\theta(x)=\tfrac12+\theta\delta,\qquad x\in D,
\]
and use the same fixed label kernel outside $D$. Denote the resulting target law by
$P_T^\theta$. The input marginal, predictors, score, and source law are fixed, so $M$ is the
same under $P_T^+$ and $P_T^-$. Lemma~\ref{lem:fibre} also gives the same law of $Z$ in the two
worlds. Consequently,
$\Law(\widetilde Z\mid P_T^+)=\Law(\widetilde Z\mid P_T^-)$ for
$\widetilde Z=(Z,M)$.

It remains to verify class membership and the benefit signs. Under $P_T^\theta$,
$\bar a_\theta=\tfrac12+\theta\delta$ and
\[
  \gamma_\theta=\theta\delta-M,
  \qquad
  |\gamma_\theta|\le\delta+|M|<\beta.
\]
Thus both laws belong to $\mathcal C_\beta$. Lemma~\ref{lem:reduction} yields
$\Delta_\theta=2\mu_T(D)\theta\delta$. The two benefits are nonzero and have opposite signs.
\end{proof}
\noindent The two evidence laws coincide exactly because only the label kernel on $D$ changes.
This is not a rate-based Le Cam argument. Their total-variation distance is zero, so the conclusion
is exact non-identifiability rather than a lower bound that varies with separation.

Randomization cannot resolve two worlds that induce exactly the same evidence law. A controller
that limits both directional errors must instead leave most of its probability on ABSTAIN.

\begin{corollary}[Matched-evidence abstention lower bound]\label{cor:matched-abstain}
Fix $\alpha\in(0,\tfrac12)$ and the shared evidence law from Theorem~\ref{lem:nonid}. Let $g$ be
any possibly randomized label-free rule. If, in every admissible world,
\begin{align*}
  \Pr[g=\adapt,\Delta\le0]&\le\alpha,\\
  \Pr[g=\freeze,\Delta\ge0]&\le\alpha,
\end{align*}
then $\Pr[g=\abstain]\ge1-2\alpha$.
\end{corollary}
\begin{proof}
The rule has the same action probabilities in both worlds. Its ADAPT probability is at most
$\alpha$ in the negative world. Its FREEZE probability is at most $\alpha$ in the positive world.
The three action probabilities sum to one.
\end{proof}

\subsection{Boundary Case}

The interior construction gives opposite nonzero signs. The boundary needs a separate argument:
one evidence-compatible world with zero benefit is already enough to defeat a uniform strict claim.

\begin{proposition}[Boundary case and closed-band abstention]\label{prop:closed-band}
Fix $\alpha\in(0,\tfrac12)$. Under Definition~\ref{def:strict-sound}, abstention is the pointwise
maximal sound action throughout $|M|\le\beta$. Any label-free rule that controls both marginal
directional errors at level $\alpha$ must abstain there with probability at least $1-2\alpha$.
Opposite nonzero benefit signs are asserted only when $|M|<\beta$.
\end{proposition}
\begin{proof}
For $|M|<\beta$, apply Theorem~\ref{lem:nonid} and
Corollary~\ref{cor:matched-abstain}.

Now suppose $|M|=\beta>0$. The declared class contains a zero-benefit world with $\gamma=-M$ and
a strict-benefit world with $\gamma=0$. These worlds have the same augmented-evidence law. In the
zero-benefit world, both directional error events occur whenever the rule commits. Dual error
control therefore gives $\Pr[\adapt]\le\alpha$ and $\Pr[\freeze]\le\alpha$. It follows that
$\Pr[\abstain]\ge1-2\alpha$, and neither strict action is uniformly sound.

Finally, if $M=\beta=0$, then $\mathcal C_0$ forces $\gamma=0$ and $\Delta\equiv0$. The same
conclusion follows directly.
\end{proof}
\noindent This convention is stricter than task regret. When $\Delta=0$, adapt and freeze are
risk-equivalent, but neither certifies a strict sign.

\subsection{Exact Strict-Commitment Frontier}

We can now answer the population question completely. The theorem separates four jobs: reduction,
sufficiency outside the band, necessity inside the band, and maximality of the resulting rule.

\begin{theorem}[Exact strict-commitment frontier]\label{thm:headline}\label{thm:frontier}
Under Assumption~\ref{ass:deploy}, fix $(\mu_T,P_S,f_0,f_a,s)$ so that $M$ is determined. Let
$\beta\ge0$. Necessity and maximality use the full, rich class $\mathcal C_\beta$. In particular,
that class contains the evidence-identical kernels from Theorem~\ref{lem:nonid} and the zero-benefit
boundary law. For a narrower subclass not closed under these constructions, only sufficiency
follows without an additional richness assumption.
\begin{enumerate}
\item[\textnormal{(i)}] \textbf{Reduction.} For every $P_T$, $\sign\Delta=\sign(M+\gamma)$
(Lemma~\ref{lem:reduction}).
\item[\textnormal{(ii)}] \textbf{Sufficiency.} If $|M|>\beta$ and $P_T\in\mathcal{C}_\beta$, then
$\sign\Delta=\sign M$.
\item[\textnormal{(iii)}] \textbf{Necessity.} If $|M|<\beta$, two laws in $\mathcal C_\beta$ have
the same augmented-evidence law and opposite nonzero benefit signs. If $|M|=\beta>0$, the
admissible residual $\gamma=-\sign(M)\beta$ gives $\Delta=0$. If $M=\beta=0$, then
$\Delta\equiv0$ throughout $\mathcal C_0$. Thus $|M|\le\beta$ supports no uniform strict direction.
\item[\textnormal{(iv)}] \textbf{Strict-commitment iff.} At the fixed augmented-evidence law, a
strict action is uniformly directionally sound if and only if $|M|>\beta$. On the closed band,
abstention is the pointwise maximal sound action in the sense of
Definition~\ref{def:strict-sound}.
\end{enumerate}
The population rule $M>\beta\Rightarrow\adapt$, $M<-\beta\Rightarrow\freeze$,
$|M|\le\beta\Rightarrow\abstain$ is the pointwise maximal sound certificate on
$\mathcal{C}_\beta$.
\end{theorem}
\begin{proof}
Part~(i) is Lemma~\ref{lem:reduction}. For~(ii), suppose $M>\beta$ and
$|\gamma|\le\beta$. Then $M+\gamma>0$, so Lemma~\ref{lem:reduction} gives $\Delta>0$. The negative
case is symmetric.

For~(iii), Theorem~\ref{lem:nonid} supplies opposite nonzero signs when $|M|<\beta$. At
$|M|=\beta>0$, the boundary law with $\gamma=-\sign(M)\beta$ has zero benefit. When
$M=\beta=0$, every law in $\mathcal C_0$ has zero benefit. These cases preclude a uniform strict
direction throughout the closed band.

Part~(iv) combines (ii), (iii), and Proposition~\ref{prop:closed-band}. Sufficiency gives the unique
sound strict direction outside the band. The impossibility and boundary results rule out both
strict directions on the band.
\end{proof}
\noindent Outside the band, the permitted residual is too small to reverse the visible direction
$M$. Inside the band, the same evidence is compatible with a failed strict claim. Abstention is
therefore part of the identified answer, not a defect of the decision rule.

\noindent The same frontier is visible from the identified benefit set. Let $d=\mu_T(D)>0$. With
the observables fixed, the full declared class gives
\[
 \Delta(\mathcal C_\beta)
 =2d\left[\max\{-\tfrac12,M-\beta\},
          \min\{\tfrac12,M+\beta\}\right].
\]
The interval inside brackets is
$[-\tfrac12,\tfrac12]\cap[M-\beta,M+\beta]$. The first set enforces
$\bar a\in[0,1]$; the second enforces $|\bar a-\tfrac12-M|\le\beta$. Every point in their
intersection is attained by a constant correctness kernel on $D$. A strict direction is available
exactly when the entire benefit interval lies above or below zero.


\begin{figure*}[t]
  \centering
  \kbgraphics[width=0.96\textwidth]{figures/fig_frontier_schematic.pdf}
  \caption{Population strict-commitment frontier. $M>\beta$ supports adapt;
  $M<-\beta$ supports freeze; and $|M|\le\beta$ supports no uniform strict commitment. Interior
  ambiguity uses opposite nonzero benefit signs, whereas the boundary uses zero-benefit ambiguity.
  This population rule is not the numerical KGA rule used on the real-data panels.}
  \label{fig:compact-frontier}
\end{figure*}

\begin{center}\footnotesize
\begin{tabular}{@{}lll@{}}
\toprule
Region & Meaning & Action \\
\midrule
$M>\beta$ & strictly positive benefit & \adapt \\
$M<-\beta$ & strictly negative benefit & \freeze \\
$|M|<\beta$ & both strict signs are feasible & \abstain \\
$|M|=\beta>0$ & zero benefit is feasible & \abstain \\
$M=\beta=0$ & benefit is identically zero & \abstain \\
\bottomrule
\end{tabular}
\end{center}

\subsection{Interpretation and Limits of the Frontier}

The theorem depends on a declared class. Choosing $\beta=0$ assumes a zero target
disagreement-conditional residual; it is not a conservative default. Without a credible upper bound on $|\gamma|$, the
population frontier cannot support a robust strict action. More unlabeled samples do not solve this
problem because the missing difference lies in the target labels, not in the input batch.

Empirical abstention is not automatically evidence of structural unknowability. A wide finite-sample
interval may instead reflect limited calibration data, estimator error, or transfer failure. The
frontier therefore leaves one central question: can its required residual bound be learned from the
same label-free evidence?

\section{What Label-Free Evidence Cannot Learn}
\label{sec:compact-budget}

\subsection{Audit Setup and Evidence Fibres}

Let $W$ collect everything a deployment-time audit may read: the unlabeled target batch, the fixed
models and score function, label-free statistics derived from them, and any source or calibration
sample whose law is specified independently of the target label kernel. An audit is a nonnegative
random quantity $\widehat\beta=g(W,U)$, where $U$ is independent randomization.

For a declared class $\mathcal C$ and an observable law $z=\Law(W)$, define the evidence fibre
\[
 \mathcal C(z)=\{P\in\mathcal C:\Law(W\mid P)=z\}
\]
and its unresolved residual radius
\begin{equation}
 \Gamma_z(\mathcal C)=\sup_{P\in\mathcal C(z)}|\gamma(P)|.
 \label{eq:compact-fibre-radius}
\end{equation}
The fibre collects target laws that the audit cannot distinguish. Its radius is therefore the
smallest uniform residual bound that those observations could support.

\subsection{Exact Fibre-Wise Audit Floor}

\begin{theorem}[Exact label-free audit floor]
\label{thm:compact-beta-minimax}
Fix a nonempty fibre $\mathcal C(z)$ and $\delta\in[0,1)$. If
$\widehat\beta$ satisfies
$\Pr_P(|\gamma(P)|\le\widehat\beta)\ge1-\delta$ for every
$P\in\mathcal C(z)$, then
\[
 \Pr\!\left(\widehat\beta\ge\Gamma_z(\mathcal C)\right)\ge1-\delta.
\]
The constant audit $\widehat\beta\equiv\Gamma_z(\mathcal C)$ is valid with $\delta=0$.
Consequently, on the high-probability event above, an audited frontier commits only where the
frontier using the declared radius $\Gamma_z(\mathcal C)$ already commits.
\end{theorem}
\begin{proof}
Every law in the fibre induces the same law for $W$, and hence the same law for
$\widehat\beta$. Validity for each fixed $P$ requires that common law to exceed
$|\gamma(P)|$ with probability at least $1-\delta$. Taking an increasing sequence of attainable
absolute residuals approaching the supremum gives the lower bound. The constant audit is valid by the
definition of the supremum. If $\widehat\beta\ge\Gamma_z$, then
$|M|>\widehat\beta$ implies $|M|>\Gamma_z$.
\end{proof}

This result does not say that useful budgets never exist. It says that the same label-free fibre
cannot create information about the target label kernel that the fibre definition removes. A valid
audit must cover the unresolved radius; reading more functions of the same $W$ does not shrink it.
The constant $\Gamma_z(\mathcal C)$ is an oracle benchmark for the fixed evidence-law fibre.
The theorem does not provide an algorithm that learns the unknown law $z$ or its radius from a
finite observed batch. We therefore describe this as a fibre-wise audit floor rather than an
unspecified minimax estimation result.

\subsection{Declared Budget and Abstention Radius}

For the full declared class
$\mathcal C_\beta=\{P:|\gamma(P)|\le\beta\}$ with
$0\le\beta\le\tfrac12$, the label-kernel construction in
Lemma~\ref{lem:fibre} gives
\begin{equation}
 \Gamma_z(\mathcal C_\beta)=\beta.
 \label{eq:compact-beta-radius}
\end{equation}
Thus for any $\beta \in [0, 1/2]$, the least fibre-wide budget a label-free audit can certify is the budget already declared by
the class. The same number is also the population abstention radius:
\[
 \boxed{\begin{gathered}
 \text{least valid fibre-wide residual bound}\\
 =\text{population abstention radius}=\beta \quad (0 \le \beta \le \tfrac12)
 \end{gathered}}
\]
This equality connects the audit problem to the strict-commitment frontier. Crucially,
Theorem~\ref{thm:compact-beta-minimax} establishes that $\beta$ must be treated as an
\emph{exogenous sensitivity parameter}, by analogy to Rosenbaum's sensitivity
parameter $\Gamma$ for unobserved confounding in observational studies~\cite{rosenbaum2002observational}
or the ambiguity set radius in distributionally robust optimization (an analogy, not a mathematical
equivalence). Extended exploratory formulations of break-even sensitivity frontiers and continuous selection odds-ratio curves are preserved in the supplementary research archive; the primary manuscript retains the exact population frontier $\Delta(\mathcal{C}_\beta)$ established in Theorem~\ref{thm:headline}.

\subsection{Why Labels at the Wrong Domain Buy Nothing}
\label{sec:compact-wrong-domain}

Suppose the audit also observes a fully labeled calibration domain, but the declared class contains
no coupling between that domain's label kernel and the deployment label kernel. The construction in
Lemma~\ref{lem:fibre} changes only deployment labels on $D$ and leaves the calibration sample
unchanged. The augmented evidence law is therefore identical across the same target fibre, so
Theorem~\ref{thm:compact-beta-minimax} applies without a smaller radius. Labels help only when they
come from the deployment domain or when a stated assumption links the labeled and deployment
domains.

\subsection{What Additional Structure Can Buy}
\label{sec:compact-additional-structure}

The missing information can be supplied, but each route changes the problem. Deployment labels can
estimate target correctness directly. Historical labeled deployments can support an
episode-marginal budget under an exchangeability model. A structural coupling can bound the change
between calibration and target domains. Domain knowledge can declare a physically meaningful
limit.

Multiclass extensions under entropy-constrained support conditions require separate assumptions and proofs; they are preserved as unproved exploratory research notes outside the maintained claims of this manuscript.

\subsection{Implications for K-Bound and KGA}

The evidence-fibre result motivates a practical question: what additional information can support
an update decision? KGA uses labeled historical shifts to predict and calibrate measured benefit.
Section~\ref{sec:compact-bridge} states its inferential target and the extra assumptions required
for population interpretation.

\section{KGA: An Empirical Benefit-Interval Gate}
\label{sec:compact-kga}

\subsection{From Population Identification to Empirical Decisions}
\label{sec:compact-bridge}

K-Bound characterizes population benefit over a declared class of target laws. KGA predicts
measured evaluation-cell benefit using labeled historical outcomes and a separately calibrated
residual interval. It does not numerically apply $|M|>\beta$, and its empirical radius
$\varepsilon$ is not an estimate of the population residual bound $\beta$. When a credible
population budget is supplied externally, a direct finite-sample frontier rule is available
(\KBSuppRef{app:direct-sampled-frontier}{Appendix}{the direct sampled-frontier construction}).
The empirical route instead depends on support for its calibration assumptions.

The binary population model gives
\begin{equation}
 \Delta=2\mu_T(D)(M+\gamma),
\label{eq:compact-bridge}
\end{equation}
while the observed cell outcome and fitted empirical prediction are
\begin{equation}
 \Delta^{\mathrm{cell}}=S(f_a;\mathcal E)-S(f_0;\mathcal E), \qquad \widehat\Delta=h_\theta(Z).
\end{equation}
Table~\ref{tab:theory-algorithm-bridge} summarizes this architectural correspondence.

\begin{table}[t]
\centering
\scriptsize
\caption{Correspondence between population theory (K-Bound) and the empirical decision rule (KGA) addressing the same paired deployment choice. Distinct inferential targets do not make the quantities equal.}
\label{tab:theory-algorithm-bridge}
\resizebox{\columnwidth}{!}{%
\begin{tabular}{@{}p{0.14\columnwidth}p{0.46\columnwidth}p{0.40\columnwidth}@{}}
\toprule
Quantity & Role in Population Frontier (K-Bound) & Status in Empirical Rule (KGA) \\
\midrule
$M$ & Observable margin on disagreement region & Not explicitly estimated \\
$\gamma$ & Hidden calibration residual ($M+\gamma=\bar a-\frac12$) & Unobserved / unestimated \\
$\beta$ & Externally declared hidden-residual bound & Not a numerical KGA input \\
$\Delta$ & Population benefit $2\mu_T(D)(M+\gamma)$ & Target estimand under population transfer \\
$\Delta^{\mathrm{cell}}$ & Not present in pure population model & Measured evaluation-cell benefit $S(f_a;\mathcal E)-S(f_0;\mathcal E)$ \\
$\widehat\Delta$ & Not present in pure population model & Fitted prediction $h_\theta(Z)$ of empirical target $\Delta^{\mathrm{cell}}$ \\
$\varepsilon$ & Not present in pure population model & Empirical residual-calibration radius \\
$b(m,\delta)$ & Not present in pure population model & Separately justified sampling allowance \\
\bottomrule
\end{tabular}%
}
\end{table}

A population interpretation requires both cell-level coverage and a separate sampling bound:
\begin{gather*}
  \Pr\bigl(|\widehat\Delta - \Delta^{\mathrm{cell}}| \le \varepsilon\bigr)
    \ge 1 - \alpha_{\mathrm{cell}}, \\
  \Pr\bigl(|\Delta^{\mathrm{cell}} - \Delta| \le b(m,\delta)\bigr)
    \ge 1 - \delta.
\end{gather*}
The cell-coverage and fresh-sample concentration premises imply, by a union bound,
\begin{multline*}
  \Pr\bigl(|\widehat\Delta - \Delta| \le \varepsilon + b(m,\delta)\bigr) \\
  \ge 1 - \alpha_{\mathrm{cell}} - \delta.
\end{multline*}
The enlarged interval therefore covers population benefit only when both premises hold.
Proposition~\ref{thm:population-transfer} states this implication; its sampling assumptions
are not established for the benchmark panels. Table~\ref{tab:theory-algorithm-bridge}
keeps the population and empirical quantities distinct.

KGA is label-free at deployment, not during development and calibration. Historical labeled cell outcomes fit $h_\theta$ and supply residuals. At the new decision, KGA reads only $Z$ and the frozen fitted artifacts. The interval premise for the empirical target is
\begin{equation}
 \Pr(|\widehat\Delta-\Delta^{\mathrm{cell}}|\le\varepsilon)\ge1-\alpha.
 \label{eq:compact-coverage}
\end{equation}
The static interval requires its stated calibration assumptions.
Sequential monitoring is a separate extension requiring its own
null model and validity conditions. The empirical results reported
here do not establish an anytime guarantee or a guaranteed response
to every harmful shift.

\subsection{The Coverage-to-Action Proposition}

The following standard implication makes the decision semantics explicit. It is not a new
calibration theorem: coverage is the premise, and establishing it is the statistical burden.

% Finite-sample certificate: split out of theory_core_main.tex during the 2026-07 restructure.
% Consumed by the "what remains" method section (sec:method), not by the impossibility section.
\begin{proposition}[Coverage-to-action implication]\label{thm:certificate}\label{thm:cert}
Fix the scalar benefit target and its evaluation unit before calibration. Let $B$ and $\widehat B$
be finite real random variables, and let $\varepsilon\in[0,\infty]$ be a measurable random radius
on a common probability space containing all calibration, benefit-model fitting,
candidate-construction, and evaluation-unit randomness. The prediction $\widehat B$ and radius
$\varepsilon$ must be available without reading the new unit's evaluation labels. For
$\alpha\in(0,1)$, suppose the interval satisfies marginal coverage
\[
\Pr\!\bigl[\,|\widehat B-B|\le\varepsilon\,\bigr]\ \ge\ 1-\alpha .
\]
Define
\[
g=\begin{cases}
\adapt & \widehat B-\varepsilon>0,\\
\freeze & \widehat B+\varepsilon<0,\\
\abstain & \text{otherwise}.
\end{cases}
\]
Then, marginally over that common probability space,
\begin{align*}
\Pr\bigl[g=\adapt,B\le0\bigr]&\le\alpha,\\
\Pr\bigl[g=\freeze,B\ge0\bigr]&\le\alpha.
\end{align*}
\end{proposition}
\begin{proof}
On the event $\{|\widehat B-B|\le\varepsilon\}$, the condition
$\widehat B-\varepsilon>0$ implies $B>0$, so a false adapt requires the complement of
this event, which has probability at most $\alpha$. The freeze case is symmetric.
\end{proof}

\begin{remark}[Unit scaling: empirical radius versus theoretical budget]\label{rem:unit-scaling}
The empirical radius $\varepsilon$ and theoretical residual bound $\beta$ live on fundamentally
different scales. The population identity is $\Delta=2d(M+\gamma)$ where $d=P_T(D)>0$ is the
disagreement probability and $2d$ is twice that probability. Thus,
$|\Delta - 2dM| = 2d|\gamma| \le 2d\beta$. The unclipped population envelope centered at $2dM$
has radius $2d\beta$ on the benefit scale, not $\beta$; the exact population benefit set
$\Delta(\mathcal C_\beta)$ in Theorem~\ref{thm:frontier} includes clipping to $[-d, d]$:
$\Delta(\mathcal C_\beta) = 2d[\max(-\tfrac12, M-\beta), \min(\tfrac12, M+\beta)]$.
Equating $\varepsilon = \beta$ is therefore a conceptual category error: it conflates a conditional
score residual on the disagreement region with an empirical cell-benefit prediction error.
While numerical equality could occur by coincidence, their mathematical definitions remain distinct.
\end{remark}

\begin{proposition}[Cell-to-population transfer via compound bound]\label{thm:population-transfer}
Let $\Delta^{\mathrm{cell}}=\frac{1}{m}\sum_{j=1}^m W_j$ be the measured cell benefit on $m$ fresh,
conditionally i.i.d.\ evaluation draws, where
$W_j=\mathbf 1\{f_a(X_j)=Y_j\}-\mathbf 1\{f_0(X_j)=Y_j\}\in[-1,1]$, with conditional mean $\Delta$.
Suppose an independently fitted benefit predictor $h$ and split-conformal radius $\varepsilon$
achieve marginal cell coverage $\Pr(|\widehat\Delta-\Delta^{\mathrm{cell}}|\le\varepsilon)\ge 1-\alpha_{\mathrm{cell}}$
over exchangeable deployment episodes. For a sampling budget $\delta\in(0,1)$, define
$b(m,\delta)=\min\{2, \sqrt{2\log(2/\delta)/m}\}$. Then the compound interval
$I_{\mathrm{population}}=[\widehat\Delta-(\varepsilon+b),\; \widehat\Delta+(\varepsilon+b)]$ satisfies
\[
\Pr\bigl(\Delta\in I_{\mathrm{population}}\bigr)\ge 1-\alpha_{\mathrm{cell}}-\delta.
\]
The corresponding population decision rule $g_{\mathrm{pop}}$ using radius $\varepsilon+b$ satisfies
\begin{align*}
\Pr\bigl(\{g_{\mathrm{pop}}=\adapt, \Delta\le0\} &\cup {} \\
\{g_{\mathrm{pop}}=\freeze, \Delta\ge0\}\bigr) &\le \alpha_{\mathrm{cell}}+\delta.
\end{align*}
\end{proposition}
\begin{proof}
By Hoeffding's inequality for bounded variables in $[-1,1]$ of range 2,
$\Pr(|\Delta^{\mathrm{cell}}-\Delta|\le b)\ge 1-\delta$. On the intersection of cell coverage
and concentration, the triangle inequality yields $|\widehat\Delta-\Delta|\le |\widehat\Delta-\Delta^{\mathrm{cell}}|+|\Delta^{\mathrm{cell}}-\Delta|\le\varepsilon+b$.
A union bound proves coverage at level $1-\alpha_{\mathrm{cell}}-\delta$. Event containment
(as in Proposition~\ref{thm:certificate}) immediately bounds the probability of an incorrect strict
direction by $\alpha_{\mathrm{cell}}+\delta$.
\end{proof}

\begin{remark}[Direct empirical sampling of the population frontier]\label{rem:direct-sampled-frontier}
An alternative route directly operationalizes the population frontier $|M|>\beta$. Given fresh
unlabeled target inputs, let $n_D>0$ fall in $D$. The empirical disagreement margin
$\widehat M=\frac{1}{n_D}\sum_{X_i\in D}s(X_i)-\frac{1}{2}$ concentrates with radius
$r_M=\sqrt{\log(2/\alpha)/(2n_D)}$. Under the externally declared bound $|\gamma|\le\beta$, the rule
$\widehat M-r_M>\beta\Rightarrow\adapt$, $\widehat M+r_M<-\beta\Rightarrow\freeze$, and $\abstain$
otherwise, guarantees directional error at most $\alpha$. This couples sampling uncertainty $r_M$
additively with structural uncertainty $\beta$.
\end{remark}

\begin{remark}[Marginal vs.\ conditional error]\label{rem:fa-marginal}
Proposition~\ref{thm:certificate} bounds the \emph{marginal, unconditional} false-commit
probabilities for the declared target $B$. Its false-adapt rate is
$\mathrm{FA}_{\mathrm u}=\Pr[g=\adapt,B\le0]$. It does \emph{not} bound
$\mathrm{FA}_{\mathrm c}=\Pr[B\le0\mid g=\adapt]$, defined when
$\Pr[g=\adapt]>0$ and reported descriptively for $B=\Delta^{\mathrm{cell}}$.
Selection-conditional coverage requires a separate construction, such as that of
Jin and Ren~\cite{jin2025selection}; conditional guarantees over specified shift classes
also require their own calibration procedure~\cite{gibbs2025conditional}.
Neither is supplied by the marginal premise above. Simultaneous protection across candidates
and time-uniform protection likewise require additional guarantees.
Risk alignment (Definition~\ref{def:risk-align}) concerns population sign identifiability;
coverage for $B$ alone yields the \emph{interval} implication above.
\end{remark}

\begin{remark}[Fallback when assumptions are unsupported]\label{cor:abstain-valid}
If interval coverage for the declared target and evaluation unit---or the assumptions supporting
it---cannot be justified, Proposition~\ref{thm:certificate} supplies no level-$\alpha$ protection
for either strict action. A separately declared operational fallback may retain the frozen model
and record abstention. Serving the frozen model without a valid negative interval is not a
certified $\freeze$ decision. This policy recommendation does not establish that the retained
model is safe or optimal. If \emph{risk alignment} (Definition~\ref{def:risk-align}) cannot
be justified for an intended population sign claim, that interpretation remains unsupported.
The interval implication of Proposition~\ref{thm:certificate} still holds when coverage for the
declared $B$ is established (cf.\ Theorem~\ref{lem:nonid}).
\end{remark}


Proposition~\ref{thm:population-transfer} formalizes the conditional bridge to population risk $\Delta$: under independent fresh evaluation draws, adding the Hoeffding concentration radius $b(m,\delta)$ to the empirical radius $\varepsilon$ yields a compound population interval with radius $\varepsilon+b$ and error budget $\alpha_{\mathrm{cell}}+\delta$. The full protocol, counterexample, and proof are detailed in \KBSuppRef{app:population-transfer}{Appendix}{the population-transfer supplement}. No such sampling-error bound is established for the reported panels; in particular, when a candidate adapts to evaluation inputs, transductive normalization couples the candidate across batches.

The formalized action rule proves that a strict decision has the correct sign whenever the
compound interval contains population benefit. This implication does not establish the
sampling or calibration premises for the benchmarks; the scope and theorem mapping are given
in \KBSuppRef{app:compact-formal}{Appendix}{the formalization appendix}.

\paragraph{ImageNet-C 108-cell population replay.}
To audit this procedure on existing benchmark data without changing historical estimands, we replayed the outcome-blind compound interval $[\widehat\Delta \pm (\varepsilon + b)]$ across 108 ImageNet-C development cells (36 cells each for Tent, EATA, and SAR; five-fold cell-outcome-disjoint cross-fitting; sensitivity parameter $m=2{,}000$, yielding uniform concentration radius $b=0.0607$; Table~\ref{tab:main-pop-replay-summary}). Full fold breakdowns and radius distributions appear in \KBSuppRef{tab:imagenetc-108-population-replay}{Table}{the supplementary 108-cell replay table} (\KBSuppRef{app:implementation-boundaries}{Appendix}{the supplementary implementation-boundary discussion}).
Finite radii were obtained on $108/108$ cells. Across the panel, the cell-level gate ($\widehat\Delta \pm \varepsilon$) issued $40$ \adapt, $11$ \freeze, and $57$ \abstain decisions with zero false adaptations ($\mathrm{FA}_u = 0.0000$) and mean oracle regret $0.0167$. Incorporating the sampling radius ($b=0.0607$, with mean compound half-width $\bar\varepsilon + b = 0.1272$) shifted compound decisions to $1$ \adapt, $10$ \freeze, and $97$ \abstain, preserving zero observed false adaptations ($\mathrm{FA}_u = 0.0000$) while shifting regret to $0.0374$. This concretely demonstrates the sensitivity to the added concentration radius $b=0.0607$: because most empirical gains on this benchmark are smaller than the universal Hoeffding penalty for $m=2{,}000$, the compound rule withholds action on nearly all adaptations, illustrating the numerical Price of Rigor rather than establishing prospective population truth. As clarified in \KBSuppRef{app:implementation-boundaries}{Appendix}{the supplementary implementation-boundary discussion}, $m=2{,}000$ is an evaluation sensitivity parameter, not fresh sampling; population truth remains unestablished on historical development records.

\begin{table}[t]
\centering
\scriptsize
\setlength{\tabcolsep}{2.5pt}
\caption{108-cell ImageNet-C population replay: cell gate ($\widehat\Delta \pm \varepsilon$) vs.\ compound population interval ($\widehat\Delta \pm (\varepsilon+b)$) under 5-fold cross-fitting ($\alpha_{\mathrm{cell}}=\delta=0.05, \alpha_{\mathrm{pop}}=0.10, m=2{,}000, b=0.0607$). Full 12-column breakdown with fold details and radii ranges in \KBSuppRef{tab:imagenetc-108-population-replay}{Table}{the supplementary 108-cell replay table}.}
\label{tab:main-pop-replay-summary}
\resizebox{\columnwidth}{!}{%
\begin{tabular}{@{}lccccc@{}}
\toprule
Candidate & $N$ & Regret (Cell / Pop) & Cell Gate (A / F / U) & Compound Pop (A / F / U) & $\mathrm{FA}_u$ \\
\midrule
\textbf{Tent} & 36 & $0.0162$ / $0.0468$ & $18$ / $0$ / $18$ & $1$ / $0$ / $35$ & $0.000$ \\
\textbf{EATA} & 36 & $0.0095$ / $0.0386$ & $20$ / $0$ / $16$ & $0$ / $0$ / $36$ & $0.000$ \\
\textbf{SAR}  & 36 & $0.0245$ / $0.0267$ & $2$ / $11$ / $23$ & $0$ / $10$ / $26$ & $0.000$ \\
\midrule
\textbf{Total / Panel} & \textbf{108} & \textbf{0.0167} / \textbf{0.0374} & \textbf{40} / \textbf{11} / \textbf{57} & \textbf{1} / \textbf{10} / \textbf{97} & \textbf{0.000} \\
\bottomrule
\end{tabular}%
}
\end{table}

\paragraph{Non-oracle population diagnostic.}
While the synthetic diagnostic in \KBSuppRef{tab:synthetic-pop-diagnostic}{Table}{the supplementary synthetic population table} tested exact Bayes prediction ($\widehat\Delta(Z) = 0.6Z$), evaluating an imperfect predictor subject to estimation noise requires an extended non-oracle study. We conducted a diagnostic across 1,000 independent trials per configuration over 16 $(N_{\mathrm{cal}}, m)$ pairs, fitting an imperfect linear predictor on finite development episodes ($m=32$; \KBSuppRef{tab:nonoracle-pop-diagnostic}{Table}{the supplementary non-oracle population table}, \KBSuppRef{app:implementation-boundaries}{Appendix}{the supplementary implementation-boundary discussion}).
An unhedged point rule ($\operatorname{sign}(\widehat\Delta)$) incurs a $30.9\%$--$35.6\%$ directional error rate ($\{g = \text{ADAPT} \wedge \Delta \le 0\} \cup \{g = \text{FREEZE} \wedge \Delta \ge 0\}$) with $0\%$ abstention (where $\mathrm{FA}_u$ is defined strictly for false \adapt; here directional errors combine false \adapt and false \freeze on null/tie states); crucially, $100\%$ of these errors occur on null/tie states ($Z=0, \Delta=0$) due to estimation noise around zero, with zero errors on strictly harmful states ($\Delta = -0.6$). Conformal cell calibration ($\widehat\Delta \pm \varepsilon_{\mathrm{cell}}$) eliminates directional error ($0.000$) with empirical cell coverage matching nominal $1-\alpha_{\mathrm{cell}}=0.95$ ($93.9\%$--$96.5\%$), and already achieves $\approx 100\%$ population coverage of $\Delta$ across all configurations in this model. The compound population interval $[\widehat\Delta \pm (\varepsilon_{\mathrm{cell}} + b)]$ likewise achieves $100\%$ empirical population coverage and zero observed wrong-direction error across these simulated trials ($0.000 \le \alpha_{\mathrm{pop}}=0.10$; the population-interval theorem bounds the population miss probability; this simulation reports zero empirically observed misses on this constructed diagnostic), while enforcing conservative abstention ($99.9\%$--$100\%$ at $m=32$, contracting safely to $\approx 33\%$ matching the null frequency for $m \ge 128$) as a direct consequence of the universal Hoeffding penalty ($b=0.4802$), rather than because cell-level calibration failed in this simulated environment.


\subsection{Evidence, Calibration, and the Action Rule}

The controlled implementation stores 11 label-free features, with algebraic rank at most 9 because two
features are deterministic differences. They summarize confidence, class balance, model
disagreement, and update size. For each dataset and adapter, a separate
\texttt{GradientBoostingRegressor}~\cite{pedregosa2011sklearn} uses squared-error loss, 250 trees,
depth 2, learning rate $0.05$, subsample $0.8$, and seed 0.
Natural protocols retain their own fitted model and feature schema; for example, CCT-20 uses a
ridge model. Details are in \KBSuppRef{app:compact-evidence}{Appendix}{the supplementary evidence specification}. A common interface does not imply
that an estimator or radius transfers across adapters.

Given residuals $r_i=|\widehat\Delta_i-\Delta_i^{\mathrm{cell}}|$, ordered as
$r_{(1)}\le\cdots\le r_{(n)}$, the order-statistic rule is
\begin{equation}
 k=\left\lceil(n+1)(1-\alpha)\right\rceil,\qquad
 \varepsilon=
 \begin{cases}
 r_{(k)},&k\le n,\\
 +\infty,&k>n.
 \end{cases}
 \label{eq:compact-rank}
\end{equation}
Split-conformal coverage requires joint exchangeability of the calibration and next-unit scores
under the fitting and candidate-construction protocol
~\cite{tibshirani2019conformal,barber2023beyond}. An independently fitted predictor and exchangeable
calibration/deployment units are one sufficient route. The rank counts evaluation units, not the
images inside them.

We distinguish two calibration protocols across the paper:
Protocol B is the current outcome-disjoint three-way fit/calibrate/score construction. Protocol A is the historical leave-one-cell-out replay.
\emph{Protocol A: Leave-one-cell-out empirical order-statistic calibration} (used in historical controlled grids). Benefit predictions $\widehat\Delta_i$ were cross-fitted on the remaining $n-1$ cells, and each scored residual $r_i = |\widehat\Delta_i - \Delta_i^{\mathrm{cell}}|$ was excluded from the pool used to compute its own radius $\varepsilon_i$. However, this historical leave-one-out construction is peer-dependent: for $j \neq i$, the model $h_{-j}$ was fitted on training data containing $\Delta_i^{\mathrm{cell}}$, so the scored cell's outcome can indirectly influence other residuals $r_j$ in its calibration pool through model fitting.
\emph{Protocol B: Three-way cell-outcome-disjoint cross-fitting} (used in the current primary CIFAR table, interval diagnostics, the 108-cell ImageNet-C population replay, and the \KBSuppRef{app:sar-official-controls}{Appendix}{the supplementary SAR controls} 27-cell audit). For each score fold, its entire complement is partitioned into strictly disjoint estimator-fit and residual-calibration subsets, ensuring that scored cell outcomes enter neither predictor training nor conformal calibration. Historical saved LOO predictions and radii do not determine these current CIFAR results.
Where historical stress grids retain pooled LOO residuals (Protocol A), they are reported with this peer-dependence caveat rather than claiming full outcome disjointness. Natural protocols use separated development, calibration, and target partitions where the archived design records that separation.

\begin{algorithm}[t]
\caption{KGA for one proposed update}
\label{alg:compact-kga}
\begin{algorithmic}[1]
\Require Frozen $f_0$, adapter $\mathcal A$, unlabeled $X$, frozen $h_\theta$ and calibration artifacts
\State Create $f_a\leftarrow\mathcal A(f_0,X)$ in isolated shadow state
\State Extract $Z=\phi(X,f_0,f_a)$ and validate the assessment inputs
\If{the assessment is unavailable}
  \State retain $f_0$, log the reason, and return \abstain
\EndIf
\State Obtain $\varepsilon$ from the calibrated residual rule in Eq.~\eqref{eq:compact-rank}
\State $\widehat\Delta\leftarrow h_\theta(Z)$; $L\leftarrow\widehat\Delta-\varepsilon$; $U\leftarrow\widehat\Delta+\varepsilon$
\If{$L>0$} \State commit $f_a$ and return \adapt
\ElsIf{$U<0$} \State discard $f_a$ and return \freeze
\Else \State retain $f_0$, log uncertainty, and return \abstain
\EndIf
\end{algorithmic}
\end{algorithm}

For example, prediction $0.04$ and radius $0.01$ give $[0.03,0.05]$ and ADAPT; prediction
$-0.04$ gives $[-0.05,-0.03]$ and FREEZE; prediction $0.005$ gives $[-0.005,0.015]$ and ABSTAIN.
These illustrate the rule, not measured results.

\begin{figure*}[t]
\centering
\kbgraphics[width=0.88\textwidth]{figures/fig_certificate.pdf}
\caption{Illustration of the three interval decisions for measured cell benefit. Only an interval
strictly above or below zero supports a commitment. The figure illustrates the rule; coverage
must be established for the intended evaluation unit.}
\label{fig:compact-certificate}
\end{figure*}

Candidate adaptation, frozen and candidate inference, evidence extraction, and shadow-state storage
all contribute to runtime. The tree prediction and interval comparison are inexpensive, but the
whole procedure is not one extra forward pass. Historical synchronized timing and peak-memory
evidence was not recovered; the retained literals are therefore audit records rather than performance
claims (\KBSuppRef{tab:system-latency-profile}{Table}{the supplementary system latency table}).
Workload-specific measurements reported in the supplement have a distinct scope, and comparable
end-to-end deployment latencies across the historical benchmark tracks remain unestablished.

\section{Primary Experimental Protocols}
\label{sec:compact-experiments}

\subsection{Two Primary Questions and the Evidence Hierarchy}

The controlled study asks whether the gate distinguishes useful from harmful updates.
The CCT-20 study asks whether a frozen gate avoids measured candidate harm under a natural
camera-location shift. These are different claims. We predefine the interpretation of each study
before presenting its results in Table~\ref{tab:compact-taxonomy}.

\begin{table*}[t]
\centering
\footnotesize
\caption{Experiment status and role. Primary presentation does not imply prospective confirmation.
All recorded unfavorable outcomes remain disclosed.}
\label{tab:compact-taxonomy}
\begin{tabular}{@{}p{0.21\textwidth}p{0.30\textwidth}p{0.40\textwidth}@{}}
\toprule
Evidence & Design and status & Question it can answer \\
\midrule
CIFAR-10-C: Tent, EATA, SAR & opened controlled grids; cross-fitted empirical calibration & candidate-specific routing on the evaluated cells \\
CCT-20: Tent & internally sealed protocol; five checkpoints and nine target camera locations & conservative retention under the stated score and predeclared utility-retention rule \\
So2Sat-LCZ42 v1/v2 & v1 development stop; v2 calibration-screen failure; targets unopened & whether the respective pre-target eligibility checks were met \\
Other replayable tracks & opened natural/semantic and candidate-configuration diagnostics & transfer limits and one-sided retention; full results in the supplement \\
PACS; iWildCam; historical ports; focused Office-Home route & partial replay, withheld metric, superseded policy, or invalid route & audit history only where the corresponding contract is incomplete or invalid \\
\bottomrule
\end{tabular}
\end{table*}

CIFAR-10-C~\cite{hendrycks2019cifar10c} supplies 2,160 recorded cells per candidate: five
archived ResNet-18 run-seed groups, each evaluated across 432
corruption/composition/update conditions. The current gate is refitted from the saved label-free
features $Z$ and measured cell outcomes under Protocol B, rather than reusing historical stored
predictions or radii. For this canonical evaluation, the analysis uses three-way
cell-outcome-disjoint cross-fitting: each scored fold has separate estimator-fit and
residual-calibration subsets, so its outcome enters neither its fitted prediction nor its radius.
This is still a retrospective opened panel, not an untouched environment test; historical stored
fields remain audit-only and do not determine the current decisions.

\subsection{CCT-20 Design and Utility Retention}
\label{sec:cct20-design}

The CCT-20 camera-location study~\cite{beery2018recognition} crosses
\CCTCheckpointCount{} independent source checkpoints with \CCTLocationCount{} target locations,
giving \CCTCellCount{} checkpoint--location cells. It uses
\CCTProbeImageCount{} probe images and \CCTEvaluationImageCount{} evaluation images
(\CCTTargetImageCount{} total). The primary metric is set-membership top-1 accuracy: a prediction
is correct if it matches any distinct archived category for that image.
It is not an official single-label leaderboard reproduction. The candidate uses the official Tent
implementation, one non-episodic Adam step at learning rate $10^{-3}$, and state reset per cell.

The protocol was internally sealed before source training and outcome access, not publicly
registered. Aggregate target label metadata had been inspected during candidate ranking; cell-level
outcomes remained unopened until one-shot scoring. This is outcome-blind execution, not a
label-unopened study. The source ridge fit used 10 rows from two camera locations; calibration used
45 rows aggregated to nine cis-test locations. Transfer to trans-test locations remains an assumption:
the protocol treats camera-location clusters as the population unit and assumes the nine trans-test
clusters are exchangeable with the declared target population under the locked sampling rule. This
is a design assumption, not an empirically proven fact about future deployments.

For either fixed policy $b$, define $G_b$ as mean baseline regret minus KGA regret.
Positive $G_b$ means lower regret for KGA. The \emph{utility-retention} rule is a predeclared,
limited harm-avoidance check. Using nominal pointwise 95\% bootstrap lower endpoints, it passes when
\begin{equation}
 L_{\mathrm{adapt}}>0,\qquad
 L_{\mathrm{freeze}}>\CCTSafeUtilityMargin.
 \label{eq:cct-safe-utility}
\end{equation}
The second threshold is the frozen-policy noninferiority margin, on the accuracy-proportion scale.
This rule can pass when KGA always retains $f_0$. The stronger routing criterion additionally
requires improvement over both fixed policies at the declared inference level and nontrivial
ADAPT/FREEZE exposure. The word \emph{safe} names this protocol endpoint, not a population guarantee.

\subsection{Scoring, Dependence, and Statistical Inference}
\label{sec:compact-inference-units}

All policies use the same evaluation records, and evaluation labels are joined only after actions
are fixed. The candidate TTA step can be transductive: episodic candidates update on the unlabeled
evaluation inputs, and online candidates may use evaluation-batch BatchNorm statistics.
Development/calibration separation therefore does not imply independent fresh-sample inference.

We report oracle regret, action counts, commitment rate, and false-commit frequencies. Empirical
interval inclusion and interval width are reported separately from commitment.
A paired regret contrast always means baseline regret minus KGA regret, so positive values favor
KGA. Accuracy and macro-F1 regrets are separated rather than pooled.

Bootstrap levels are nominal throughout. CIFAR sensitivity resamples six corruption families,
conditional on an archived checkpoint. CCT-20 uses a paired two-way product bootstrap over
checkpoint rows and location columns; sign-flip calculations instead use nine location means.
For a vector $G$ of cluster gaps, the sign-flip reference model is
\begin{equation}
 G\overset{d}{=}s\odot G
 \quad\text{for every }s\in\{-1,+1\}^{m}.
 \label{eq:compact-signflip-null}
\end{equation}
Independent zero-symmetric cluster gaps suffice. Exchangeability, marginal symmetry, or a single
global sign symmetry alone is not sufficient~\cite{hemerik2018exact}. Enumeration removes Monte
Carlo error but does not verify this null model. Retrospective Holm adjustment over the six prospectively named contrasts is reported, but Holm adjustment requires valid input $p$-values;
nominal family intervals also require valid marginal coverage of their constituent intervals.
The full inference specification and retrospective search status are collected in
\KBSuppRef{app:compact-inference}{Appendix}{the supplementary inference specification}.

\section{Primary Results}
\label{sec:compact-results}

\subsection{Controlled Routing Depends on the Candidate}

Table~\ref{tab:primary-cifar} reports all three CIFAR-10-C candidates under the same current rule.
Tent and EATA have lower measured point regret than always-adapt and always-freeze.
SAR is the counterexample: always-adapt has lower regret than KGA. The mixed panel therefore
supports candidate-specific routing, not a claim that gating improves every adapter.

\begin{table*}[t]
\centering
\footnotesize
\caption{Primary controlled CIFAR-10-C results. Regret is the mean accuracy shortfall from the
per-cell oracle, on the proportion scale; lower is better. A/F/U denotes ADAPT/FREEZE/ABSTAIN.
Commitment rate is $(A+F)/n$, not interval coverage. Cells share checkpoints and residual pools.}
\label{tab:primary-cifar}
\resizebox{\textwidth}{!}{% AUTO-GENERATED by scripts/make_tables.py. Display only; authorities unchanged.
\begin{tabular}{@{}lrrrrcrr@{}}
\toprule
Candidate & $n$ & KGA & Adapt & Freeze & A/F/U & $\mathrm{FA}_{\mathrm u}$ & Commitment rate \\
\midrule
Tent & 2160 & 0.0018 & 0.0080 & 0.1239 & 1094/334/732 & 0.0000 & 0.6611 \\
EATA & 2160 & 0.0016 & 0.0033 & 0.1313 & 1207/118/835 & 0.0000 & 0.6134 \\
SAR & 2160 & 0.0017 & 0.0003 & 0.1405 & 1430/0/730 & 0.0005 & 0.6620 \\
\bottomrule
\end{tabular}
}
\end{table*}

Tent regret is $\CIFARtentKga$, compared with $\CIFARtentAdapt$ for always-adapt and
$\CIFARtentFreeze$ for always-freeze. EATA's values are $\CIFAReataKga$,
$\CIFAReataAdapt$, and $\CIFAReataFreeze$. In the canonical primary panel (Table~\ref{tab:primary-cifar}), their actions are 1,094/334/732 and
1,207/118/835, respectively; both have zero observed false adaptations ($\mathrm{FA}_u = 0.0000$).
SAR commits 1,430/0/730 actions with regret $0.0017$, versus $0.0003$ for always-adapt ($\mathrm{FA}_u = 0.0005$, 1 false adaptation).

On the recorded CIFAR SAR panel, always-adapt has very low paired-oracle regret. KGA's higher regret
reflects the net benefit forgone on cells where it retains the frozen predictor. The current results
do not isolate whether those decisions arise mainly from prediction error, conservative calibration,
or candidate-update dynamics. Normalization and sample-filtering effects remain hypotheses rather than
established causes.

In contrast, Tent and EATA suffer degradation under severe corruptions, allowing KGA's selective
freezing and abstention to achieve net regret reductions over always-adapt ($+0.0062$ for Tent,
$+0.0017$ for EATA on the canonical primary panel; matching the synchronized baseline study).
Figure~\ref{fig:decisive-pareto} illustrates this regret-minimization behavior across deployment streams with varying harmful fraction $p$: while unconstrained always-adapt degrades linearly as harmful mass increases, KGA tracks the lower mixture-regret envelope, maintaining lower empirical regret across the constructed mixture than the fixed policies over most mixture proportions.

\begin{figure}[t]
  \centering
  \kbgraphics[width=\columnwidth]{figures/fig_decisive_pareto_cifar10c.pdf}
  \caption{CIFAR-10-C stress grid mixing-ratio regret envelope: mean regret vs.\ the harmful fraction $p$ of the deployment stream for Tent, comparing always-adapt, always-freeze, and KGA. While unconstrained always-adapt degrades sharply as the harmful fraction increases and always-freeze stalls on helpful distributions, KGA tracks the lower mixture-regret envelope, maintaining lower empirical regret across the constructed mixture than the fixed policies over most mixture proportions.}
  \label{fig:decisive-pareto}
\end{figure}

A retrospective cluster-bootstrap sensitivity analysis across the six observed CIFAR-10-C corruption
families (\KBSuppRef{tab:compact-current-family}{Table}{the supplementary family-sensitivity table})
indicates an estimated paired regret reduction for Tent of $+0.0062$ over always-adapt
(95\% CI $[+0.0023,\,+0.0095]$) and $+0.1221$ over always-freeze (95\% CI $[+0.0630,\,+0.1777]$).
Across five stream and evaluation runs (seeds 0--4) on the archived source checkpoint, per-seed regret reductions for Tent are $+0.0058, +0.0065, +0.0061, +0.0072, +0.0054$, strictly positive across each evaluated run. These repetitions evaluate stream variation rather than variation across independently trained source models. Historical run-seed bootstrap vectors from earlier development cycles are cataloged separately in the supplement. Corruption-family sensitivity breakdowns appear in \KBSuppRef{tab:compact-current-family}{Table}{the supplementary family-sensitivity table}.

\subsection{Observed Interval Inclusion Is a Separate Diagnostic}
\label{sec:interval-diagnostics}

Commitment rate measures how often the interval excludes zero. Interval inclusion measures how
often it contains the observed cell outcome. Table~\ref{tab:interval-diagnostics} reports these
diagnostics separately under Protocol B (disjoint fit/calibrate/score cross-fitting) at nominal target
$1-\alpha=0.90$.

\begin{table*}[t]
\centering
\footnotesize
\caption{Retrospective interval diagnostics for the controlled cells under Protocol B (disjoint fit/calibrate/score cross-fitting), at nominal target
$1-\alpha=0.90$. Observed inclusion concerns $\Delta^{\mathrm{cell}}$.
Candidate/family breakdowns and directional-error denominators are provided in the supplement.}
\label{tab:interval-diagnostics}
\resizebox{\textwidth}{!}{% Generated by build_current_policy_interval_diagnostics.py; do not edit.
% RETROSPECTIVE dependent-cell description, not independent calibration validation.
% LOO pool reuse strongly rank-constrains pooled inclusion. All rates are fractions.
% Full widths are 2*epsilon, unclipped, in accuracy proportion units. -- is undefined.
\begin{tabular}{lrrrrrrrr}
\toprule
Candidate & $n$ & Nominal & Inclusion & Mean width & Median width & Commitment & $\mathrm{FA}_{u}/\mathrm{FA}_{c}$ & $\mathrm{FF}_{u}/\mathrm{FF}_{c}$ \\
\midrule
TENT & 2160 & 0.9000 & 0.9046 & 0.0469 & 0.0483 & 0.6611 & 0.0000/0.0000 & 0.0000/0.0000 \\
EATA & 2160 & 0.9000 & 0.9060 & 0.0387 & 0.0378 & 0.6134 & 0.0000/0.0000 & 0.0005/0.0085 \\
SAR & 2160 & 0.9000 & 0.8949 & 0.0279 & 0.0277 & 0.6620 & 0.0005/0.0007 & 0.0000/-- \\
\bottomrule
\end{tabular}
}
\end{table*}

Observed inclusion rates closely track the 90\% nominal level across all three candidates:
Tent includes 1,954 of 2,160 outcomes (90.46\%), EATA includes 1,957 (90.60\%), and SAR includes 1,933 (89.49\%).
This retrospective diagnostic is rank-constrained and is not independent validation of interval coverage under unseen distributions.
Evaluating directional errors relative to committed decisions reveals asymmetric behavior:
EATA commits one false FREEZE among 118 FREEZE decisions ($\mathrm{FF}_c = 0.0085$, occurring in Gaussian noise).
Tent commits zero false FREEZE decisions among 334 FREEZE decisions ($\mathrm{FF}_c = 0.0000$), while SAR commits
zero FREEZE decisions and one false ADAPT among 1,430 ADAPT decisions ($\mathrm{FA}_c = 0.0007$, in pixelate).
These conditional denominators describe the observed directional errors among committed decisions
on this panel, reflecting low empirical error rates on the evaluated records.

\subsection{Synchronized Development Baseline Comparison}
\label{sec:synchronized-baselines}

Table~\ref{tab:synchronized-baselines} compares KGA against fixed policies, scalar heuristics, and an uncalibrated
point gate on 2,160 CIFAR-10-C Tent rows, using identical saved features and candidate-specific oracles under Protocol B.

\begin{table*}[t]
\centering
\small
\caption{Current outcome-disjoint saved-feature comparison on 2,160 CIFAR-10-C
Tent rows, computed from the same paired scores. All rows
have the same oracle and weights; both fitted gates use identical predictions.
Lower regret is better. $\mathrm{FA}_{\mathrm u}$ is the fraction of all rows
with ADAPT and $B\le0$; commitment is $(A+F)/n$.}
\label{tab:synchronized-baselines}
\resizebox{\textwidth}{!}{% Generated arithmetic-only opened-record replay; no native methods or significance claims.
\begin{tabular}{@{}lrrrrrrr@{}}
\toprule
Policy & Regret & Mean acc. & $\mathrm{FA}_{\mathrm u}$ & Commitment & ADAPT & FREEZE & ABSTAIN \\
\midrule
Always Adapt & 0.0080 & 0.7869 & 0.334 & 1.000 & 2160 & 0 & 0 \\
Always Freeze & 0.1239 & 0.6709 & 0.000 & 1.000 & 0 & 2160 & 0 \\
Confidence increase & 0.0080 & 0.7869 & 0.250 & 1.000 & 1834 & 326 & 0 \\
Entropy decrease & 0.0081 & 0.7868 & 0.249 & 1.000 & 1812 & 348 & 0 \\
Point benefit (same prediction) & 0.0004 & 0.7945 & 0.041 & 1.000 & 1471 & 689 & 0 \\
Current KGA & 0.0018 & 0.7931 & 0.000 & 0.661 & 1094 & 334 & 732 \\
\bottomrule
\end{tabular}
}
\end{table*}

KGA achieves lower point regret ($0.0018$, mean accuracy $0.7931$) than either fixed policy
(always-adapt $0.0080$, always-freeze $0.1239$), while committing zero false adaptations
($\mathrm{FA}_u = 0.000$, zero false ADAPT decisions out of 1{,}094 ADAPT commitments).
The same-prediction point gate achieves even lower point regret ($0.0004$, mean accuracy $0.7945$)
by committing to adaptation whenever $\widehat\Delta > 0$ without an uncertainty margin;
however, this unhedged stance causes 88 non-positive adaptations ($\mathrm{FA}_u = 0.041$,
$88/2{,}160$). KGA trades point regret for greater directional conservatism: by withholding commitment on 732 ambiguous cells, it records zero false ADAPT decisions on these evaluated cells.
The scalar heuristics (confidence increase and entropy decrease) achieve regret of $0.0080$ and $0.0081$,
respectively, with substantial false adaptations ($\mathrm{FA}_u = 0.250$ and $0.249$) and lack formal calibration.
An expanded nine-policy decision rule comparison—including drift monitors, ATC-style estimators, and fixed margins—is reported in \KBSuppRef{tab:cifar10c-decision-rules}{Table}{the supplementary decision rule table}.

\subsection{CCT-20: Conservative Retention, Not Selective Routing}
\label{sec:cct20-result}

On the recorded CCT-20 cells, KGA served the frozen model throughout and avoided the measured loss of always adapting.
Across \CCTCellCount{} evaluation cells, KGA issued \CCTAdaptCount{} ADAPT, \CCTFreezeCount{} FREEZE,
and \CCTAbstainCount{} ABSTAIN decisions. Because every non-adapt decision serves the base
model, KGA served frozen predictions across the entire benchmark.

\begin{table*}[t]
\centering
\footnotesize
\caption{CCT-20 safe-utility endpoint defined in Eq.~\eqref{eq:cct-safe-utility}.
Contrasts are baseline regret minus KGA regret on the set-membership accuracy-proportion scale.
Contrasts report nominal pointwise 95\% two-way bootstrap intervals.
Location-level records are detailed in the supplement.}
\label{tab:cct20-safe-utility}
% AUTO-GENERATED by scripts/make_tables.py. Sealed endpoint, display rounding only.
\begin{tabular}{@{}lrrl@{}}
\toprule
Comparator & Mean contrast & Nominal 95\% CI & Required lower bound \\
\midrule
Always adapt & 0.1819 & [0.0937, 0.2914] & $L>0.0000$ \\
Always freeze & 0.0000 & [0.0000, 0.0000] & $L>-0.0050$ \\
\bottomrule
\end{tabular}

\end{table*}

KGA matched always-freeze and avoided the measured degradation of always-adapt. The predeclared utility-retention endpoint was therefore met. The stronger selective-routing criterion was not: KGA issued no ADAPT decisions and failed to select the single helpful candidate.
All recorded development cells for Tent had negative measured benefit.

\subsection{Development Stops and Adverse Diagnostics}
\label{sec:so2sat-development-stop}

The earlier v1 study on So2Sat-LCZ42~\cite{zhu2020so2sat} tested \SoTwoCandidateCount{} candidates on
\SoTwoDevelopmentCityCount{} development cities and \SoTwoCheckpointCount{} independent
checkpoints. Tent had both helpful and harmful cities, but neither candidate passed all
predeclared eligibility checks. The study stopped before gate calibration: no target inputs or
labels were accessed, and there is no target natural-shift score. The development outcomes are
now open, so retuning them cannot produce fresh confirmation.
The separately locked v2 calibration study is reported in
Section~\ref{sec:nextphase-natural-cost}; it also stopped before target access,
after failing its fixed calibration eligibility screen.

The auxiliary studies also show limitations. ImageNet-R and PACS favor always-adapt;
ImageNet-C is sensitive to the candidate and its operating point. Office-Home and RxRx1 mainly
retain the frozen model. The opened one-sided Camelyon17 panel adapts throughout because all
18 cells are helpful.
A retrospective 72-condition Camelyon17 diagnostic contains both helpful and harmful
candidate effects but issues 0 \adapt, 2 \freeze, and 70 \abstain decisions, matching
always-freeze regret ($0.0438$; \KBSuppRef{app:camelyon17}{Appendix}{B.5}). Together with
CCT-20 and RxRx1, it provides retention evidence without ADAPT exposure. These studies do not
establish prospective selective routing on unseen natural shifts.

\subsection{Additional Tests of Decision Value}
\label{sec:nextphase-comparisons}

The following analyses ask three related questions at different levels. First, does residual
calibration add decision value beyond a shared point predictor? Second, can paired multiclass
transport make a decision under an explicitly declared structural assumption? Third, does a
natural-shift calibration screen provide enough evidence to authorize target scoring? Sections
\ref{sec:nextphase-calibration}--\ref{sec:nextphase-natural-cost} examine these questions
separately; Section~\ref{sec:nextphase-deployment} describes the local deployment integration. None turns a conditional or retrospective result into a general deployment guarantee.

\subsection{What Does Calibration Add to a Shared Predictor?}
\label{sec:nextphase-calibration}

We conducted a new retrospective comparison on the already opened CIFAR, ImageNet,
and mixed-checkpoint records. Every arm shares the same fitted GBRT within a fold.
CIFAR and ImageNet folds exclude complete scored corruption types from fitting,
tuning, and calibration. Checkpoint-only folds exclude the scored source model
and rotate held-out semantic groups; joint folds exclude both the scored source
model and the complete corruption type.
This tests transfer beyond the earlier cell-wise splits, while retaining the historical
data's provenance and dependence limits. It is not a new independent deployment study.
The loss grid, fold construction, and policy family were fixed before the new fitting;
held-out outcomes enter scoring only after decisions have been saved.

The primary loss uses a prespecified harm weight of five:
\begin{equation}
 L_5=\frac{1}{N}\sum_{i=1}^N
 \left[(1-a_i)(B_i)_+ +5a_i(-B_i)_+\right],
 \label{eq:nextphase-loss}
\end{equation}
where $B_i$ is measured paired cell benefit and $a_i$ is one for ADAPT and zero
otherwise. A fixed margin minimizes this loss
on tune-only groups. Constant and locally scaled residual intervals use the exact
conformal rank, without clipping a rank that exceeds the calibration sample size.
We report $100L_5$ in percentage points, so lower values indicate lower weighted decision loss.
The full study includes harm weights $1,5,20$, nominal levels $80,90,95\%$,
group-maximum residuals, exposure-matched controls, and binomial risk-selection arms.

\begin{table*}[t]
\centering
\small
\caption{Retrospective shared-predictor comparison. Losses are $100L_5$ in percentage
points (lower is better). Margin is selected on tuning data; Constant and Scaled
use nominal 90\% cell residual intervals. Inclusion C/S is observed held-out
interval inclusion for those two arms, not a coverage guarantee. The same 108
mixed-checkpoint records appear in two evaluation designs; they are not 216 independent observations.}
\label{tab:nextphase-calibration}
\resizebox{\textwidth}{!}{% Generated by scripts/make_next_phase_tables.py; do not edit numerical rows.
% Original authority: output/next_phase/calibration_value_v1/FINAL_RESULTS_SUMMARY.json
% Byte-identical rendering copy: next_phase/table_sources/calibration_value_summary.json
% SHA-256: 1b2f00d9dba566d35e86311f0c5e4e63833c333b702cee11c29eeb2411d9a191
\begin{tabular}{@{}lrrrrrrr@{}}
\toprule
Cohort / outer hold-out & Cells & Point & Margin & Constant & Scaled & Inclusion C & Inclusion S \\
\midrule
CIFAR EATA / corruption & 2160 & 0.1077 & 0.0663 & 0.1556 & 0.0914 & 54.35\% & 51.67\% \\
CIFAR SAR / corruption & 2160 & 0.1677 & 0.0465 & 0.1291 & 0.0844 & 50.51\% & 48.19\% \\
CIFAR Tent / corruption & 2160 & 0.2867 & 0.2264 & 0.2150 & 0.1429 & 60.69\% & 58.24\% \\
ImageNet EATA / corruption & 135 & 0.0481 & 0.4443 & 0.6939 & 1.8787 & 91.11\% & 92.59\% \\
ImageNet SAR / corruption & 135 & 0.0000 & 0.2115 & 2.5989 & 1.8754 & 97.04\% & 93.33\% \\
ImageNet Tent / corruption & 135 & 2.1913 & 1.3756 & 1.4465 & 1.4465 & 93.33\% & 96.30\% \\
Mixed Tent / checkpoint & 108 & 0.3236 & 0.4032 & 0.7954 & 5.2708 & 85.19\% & 90.74\% \\
Mixed Tent / checkpoint + corruption & 108 & 0.5829 & 0.4611 & 0.9403 & 2.6329 & 79.63\% & 87.96\% \\
\bottomrule
\end{tabular}
}
\end{table*}

Table~\ref{tab:nextphase-calibration} does not establish a general advantage for
calibration. Scaled calibration reduces CIFAR Tent loss from $0.2867$ to $0.1429$
percentage points versus point prediction, and from $0.2264$ versus the tuned
margin. Its loss at matched exposure is only modestly lower than point ranking
($0.1429$ versus $0.1610$). Simple margins do better on CIFAR EATA and SAR;
point or margin policies do better in the other cohorts. Constant-radius gates
have exactly the point predictor's ranking at equal exposure: their role is
choosing an operating point, rather than improving ranking.

Nominal 90\% constant-cell intervals include only $50.51$--$60.69\%$ of held-out
CIFAR benefits; scaling and group maxima do not restore nominal inclusion. Decision loss and
interval inclusion answer different questions: a lower observed loss does not establish nominal
coverage, and observed inclusion does not establish useful routing.
All independently held-out environment-calibration arms abstain: this split
provides no calibration environments independently excluded from fitting and
tuning. The 10\% conditional group-risk arm also abstains throughout. Even with
zero observed errors, its 13-threshold family needs 53 exposed calibration groups
at within-fold confidence error $.05$, exceeding the largest partition of 45.
These are information and transfer limits, rather than evidence that every
possible policy is unsafe. The 20\% sensitivity selects some CIFAR updates;
its held-out group error rates differ materially from cell rates
(\KBSuppRef{app:nextphase-calibration}{Appendix}{the supplementary calibration comparison}).
B6 remains descriptive: this study does not recover its missing original execution
identities or repair its historical rank clipping.

\subsection{A Conditional Paired-Transport Certificate}
\label{sec:nextphase-paired}

The population frontier leaves open how extra structural information can make a
decision possible. We instantiate one such assumption regime for fixed multiclass
predictors. Let $r=(f_0(X),f_a(X))$ index their joint prediction bins,
$A_{ry}=P_S(r\mid Y=y)$, and $q_r=P_T(r)$. With unknown target prior $\pi$, declare
\begin{equation}
 \sum_y\pi_y\TV\bigl(P_T(r\mid y),A_{\cdot y}\bigr)\le\rho.
 \label{eq:nextphase-transport}
\end{equation}
The budget $\rho$ requires external justification; unlabeled evidence cannot
generally validate it. Both predictors must be fixed independently of the
source and target evidence samples. In particular, adapting on the same target
examples later used as independent evidence does not satisfy this construction.

For $R$ bins and $K$ classes, simultaneous binomial intervals for the $RK+R$
source-conditional and target-marginal probabilities define a linear feasible
set. The variables are target joint masses $t_{ry}$, transported source masses
$s_{ry}=A_{ry}\pi_y$, priors, and absolute-value slacks for
Equation~\eqref{eq:nextphase-transport}. Source interval endpoints multiply
$\pi_y$, so the resulting constraints remain linear. Minimize and maximize
\begin{equation}
 \Delta(t)=\sum_{r,y}
 \bigl(\mathbf{1}\{f_a(r)=y\}-\mathbf{1}\{f_0(r)=y\}\bigr)t_{ry}.
 \label{eq:nextphase-paired-objective}
\end{equation}
On the simultaneous interval event, the true table is feasible and these extrema
enclose its benefit. Under the stated sampling and transport assumptions,
strict interval thresholding therefore has unconditional wrong-direction
probability at most $\alpha$. It does not give acceptance-conditional or
repeated-use error control. Detected infeasibility, solver failure, or an interval
touching zero produces \abstain. The full constraints, proof, and numerical
scope appear in \KBSuppRef{app:nextphase-paired-proof}{Appendix}{the supplementary paired-transport proof}.

This construction applies existing interval-constrained LP inference
~\cite{li2026matrixconstraints} to a paired contrast. Joint optimization can
retain information lost by separately bounding the two accuracies: in one exact
three-class construction, the paired interval is $[.1,.4]$, whereas subtracting
separate extrema on the same feasible set gives $[-.1,.6]$. An unidentifiable
prior can therefore coexist with an identified positive action. This is a
concrete use of dependence, not a first claim for action identification
~\cite{jhawar2026identifiable} or new LP machinery.

\begin{table*}[t]
\centering
\small
\caption{Selected diagnostic regimes from 4,860 synthetic trial configurations
(nine scenarios, three sample sizes, three sensitivity budgets, 60 trials each).
A/F/U count ADAPT/FREEZE/ABSTAIN over 60 trials. Separate subtracts accuracy
extrema on the same polytope; it is not an implementation of MaC-LP. The
same-evidence violation row deliberately violates the declared assumption and
reuses the paired-cancellation row's observable count streams.}
\label{tab:nextphase-paired}
\resizebox{\textwidth}{!}{% Generated by scripts/make_next_phase_tables.py; do not edit numerical rows.
% Original authority: output/next_phase/paired_transport_v1/summary.json
% Byte-identical rendering copy: next_phase/table_sources/paired_transport_summary.json
% SHA-256: e24468294bdbc1308e3def160fe83b1a56c6c306f276e068a83db8c5b4227b29
\begin{tabular}{@{}lrrrrrr@{}}
\toprule
Synthetic regime & $n$/class & $\rho$ & True $\Delta$ & Paired A/F/U & Separate A/F/U & Plug-in A/F/U \\
\midrule
Paired cancellation & 100 & 0.0 & 0.16 & 60/0/0 & 0/0/60 & 60/0/0 \\
Same evidence, violated assumption & 100 & 0.0 & -0.50 & 60/0/0 & 0/0/60 & 60/0/0 \\
Near null & 100 & 0.0 & 0.02 & 0/0/60 & 0/0/60 & 38/22/0 \\
Near null & 10000 & 0.0 & 0.02 & 20/0/40 & 0/0/60 & 60/0/0 \\
Helpful & 1000 & 0.0 & 0.20 & 60/0/0 & 60/0/0 & 60/0/0 \\
Helpful, wider sensitivity & 1000 & 0.1 & 0.20 & 0/0/60 & 0/0/60 & 60/0/0 \\
Harmful & 1000 & 0.0 & -0.20 & 0/60/0 & 0/60/0 & 0/60/0 \\
\bottomrule
\end{tabular}
}
\end{table*}

The paired procedure adapts in all 60 cancellation trials while the separate
accuracy relaxation always abstains. Near a small positive benefit, it avoids
the plug-in rule's 22 wrong FREEZE decisions out of 60 at the smaller sample
size; its price is abstention. Larger sensitivity budgets likewise reduce
commitment. Most importantly, the same observed evidence can have true benefit
$-.5$ under a violated transport assumption: the paired procedure then makes
60 harmful ADAPT decisions. The diagnostic demonstrates both conditional value
and the impossibility of obtaining the missing assumption from unchanged
unlabeled observations. It does not establish useful natural-shift adaptation.

% Final source-bound natural and deployment results are integrated here.
\subsection{Natural-Shift Calibration and Execution Costs}
\label{sec:nextphase-natural-cost}

The separate So2Sat v2 study fixes the controller derived from nine already
opened development cities and uses five independently trained source checkpoints.
Its calibration unit is a city: the score is the maximum absolute residual across
all five checkpoints, with the 18th order statistic among 19 city scores at
$\alpha=.10$. Under exchangeability, this construction targets simultaneous
coverage of the five checkpoints for one new city. It does not give simultaneous
coverage of ten target cities, coverage conditional on passing an eligibility
screen, or a demonstrated guarantee under the held-out geographic shift.
The recorded chronology is an internal execution record rather than external
proof of universal historical nonaccess.

The locked calibration completed all 95 city--checkpoint cells. Its radius was
$0.14224$ (14.224 percentage points), and its decisions were 2 \adapt,
0 direct \freeze, and 93 \abstain. Both accepted updates were helpful in the
calibration records, but they occurred in only one city; no city supplied a
direct \freeze. Thus both prespecified requirements of at least seven direct-action
cities failed, and the study stopped with zero target pixels or labels read.
No target authorization or target score was produced. The observed mean calibration
utility was $0.0989$ percentage points for the gate, $0.1939$ for always-adapt,
$0.1264$ for the checkpoint-only point baseline, and zero for always-freeze.
These same calibration outcomes set the radius and screen, so these utilities
are descriptive and cannot demonstrate independent target benefit.

A posthoc diagnosis of those opened records found 51 helpful, 43 harmful,
and one tied candidate effect. The radius was set by S\~ao Paulo checkpoint~0:
its predicted benefit was 14.4215 percentage points but its observed benefit
was only 0.1976 points. The resulting radius exceeded every negative prediction
magnitude (at most 2.6067 points), preventing direct \freeze. Thus the failed
screen reflects insufficient prediction accuracy and interval informativeness,
rather than an absence of helpful and harmful candidates. This diagnosis did
not refit the controller, alter the screen, or authorize target access.

Calibration took 36.55 minutes including input verification and maintenance
overhead. The 95 full cell executions summed to 33.63 minutes; median cell cost
was 15.15 seconds, p95 was 63.34 seconds, and the maximum was 65.65 seconds.
Process peak RSS was 697 MiB, and the largest MPS driver allocation sampled after
a cell was 1,164 MiB. The latter is not a continuously measured peak.

The controller averages probe features across all five checkpoints. It therefore
pays for five shadow candidates before routing a city, including when its final
action retains the frozen model. Calibration maintenance, candidate construction,
probe inference, and subsequent serving inference must be counted separately.

A separate stage profile used the already-opened Cairo development city,
with 8,606 probe images and 8,839 evaluation images, all five checkpoints,
and three repetitions on this host's MPS device. The sum of five checkpoint
resets, source probes, Tent updates, candidate probes, feature extraction,
and city-mean prediction was 203.22--208.24 seconds per repetition
(median 207.44 seconds), excluding subsequent serving inference and cleanup.
Across 15 checkpoint executions, median source-probe/statistic computation,
Tent update, and candidate-probe inference took 24.33, 9.66, and 6.89 seconds.
Inference on the evaluation images took a median 7.04 seconds for a source
model or 7.06 seconds for a candidate; both alternatives were measured separately,
not served as an ensemble. Setup and identity verification took 69.10 seconds;
the entire profile took 902.84 seconds. Process peak RSS was 675 MiB and the
largest sampled post-stage MPS driver allocation was 1,138 MiB.
These are descriptive costs from one opened city, without label reads,
explicit warmup, or exclusive-host isolation; document builds overlapped
some repetitions. They establish neither new efficacy nor a production latency
guarantee. Calibration acquisition, energy, and long-duration maintenance
costs remain unmeasured.

\subsection{A Local Deployment Integration}
\label{sec:nextphase-deployment}

The reference lifecycle keeps source checkpoint bytes immutable and binds each
decision to source, candidate, input batch, adapter, estimator, calibration,
and schema identities. It rechecks calibration expiry after candidate execution
and journal persistence, rejects malformed or outcome-bearing evidence, and
preserves the source on failed estimates or journal writes. A validity window
and engineering request limits do not establish statistical validity under
adaptive repeated use. Callbacks remain trusted application code; content hashes
authenticate identities, not the truth of calibration or transfer assumptions.
The contract also caps retained receipts and journal bytes; capacity exhaustion
permanently closes the session and requires source fallback. Model retention is
bounded by the configured assessment count and must be sized to checkpoint bytes.
In 1,000 measured byte-fixture requests, median lifecycle overhead was
$0.114$\,ms and p95 was $0.184$\,ms; all 1,027 journal records, including warmups
and fault cases, passed hash-chain verification. These figures exclude neural
candidate execution. Forty-six lifecycle regressions cover expiry, failed writes,
concurrent retries, invalid evidence, and capacity boundaries.

We also exercise the existing certificate API in local Docker production mode,
with a read-only root filesystem, a non-root process, authentication, and a
loopback-only port. This endpoint processes supplied certificate inputs; it
does not execute neural candidates or integrate the new lifecycle automatically.
These are local integration and load measurements, not independent production
deployments or a field reliability result.
For 100 sequential requests, median latency was $2.40$\,ms and p95 was
$3.78$\,ms; with four concurrent workers and 100 requests, these were
$6.00$\,ms and $12.90$\,ms. The image's 54 code and lock-file hashes matched
the current source. This short local load test does not measure long-duration
availability, energy, or independently replicated production utility.


\section{Limitations and Broader Impact}
\label{sec:compact-sensitivity}

\subsection{Assumptions, Not an Unconditional Safety Guarantee}

The population frontier is exact within its declared rich binary class. A credible residual bound
must come from domain knowledge, target labels, or a stated transfer assumption. A narrower class
may permit stronger decisions, while an unsupported bound cannot justify a strict population claim.
The useful object is the class-dependent feasible set, not the algebraic inequality alone.

\paragraph{Justifying the population residual bound.}
\label{par:beta-surrogate-failure}
A development-domain proxy does not resolve the population residual bound required by
Theorem~\ref{thm:compact-beta-minimax}. Within the declared class, $\beta$ must come from
external domain knowledge, target labels, or an explicit transfer assumption.
Processing the same label-free evidence cannot supply that missing restriction.
The supplementary discussion records the withdrawn proxy analysis
(\KBSuppRef{app:withdrawn-beta-proxy}{Appendix}{the withdrawn-proxy discussion}).

For KGA, empirical interval coverage requires separate support at the intended decision unit.
A larger radius reduces commitment but can increase regret when the frozen model is worse.
A smaller radius increases exposure and makes calibration failures more consequential.
Changing the adapter, checkpoint, feature schema, or update settings changes the prediction problem
and requires recalibration or a justified transfer argument.
The new held-out-corruption comparison directly exhibits this failure: nominal
90\% calibration transfers poorly on CIFAR, and the best useful policy depends
on the candidate and harm weight. The paired-transport certificate supplies a
different sufficient assumption regime; it does not validate either that
calibration transfer or its own transport budget from unlabeled target data.

\paragraph{Repeated deployment needs a separate guarantee.}
A predeclared policy does not by itself supply exchangeability or time-uniform safety.
Reusing calibration data, changing the frozen baseline after accepted updates, and selecting later
candidates from earlier outcomes change the joint process. Adaptive conformal inference studies
online coverage frequency with feedback~\cite{gibbs2021adaptive}; that is not time-uniform control
of KGA's false commitments. Selection-conditional and group-conditional guarantees likewise require
additional methods, as discussed in Section~\ref{sec:compact-related}.

\subsection{What the Experiments Establish---and What Remains Open}
\label{sec:empirical-delimitations}

Four limitations determine how these experiments should be interpreted:
\begin{enumerate}
  \item \textbf{Retention vs.\ Opportunistic Routing under Natural Shift:} On real-camera distribution shifts
  (CCT-20), KGA served the frozen model across all 45 cells and avoided the measured degradation of always-adapt
  in that study. However, it made zero \adapt decisions. Prospectively demonstrating selective opportunistic routing on an unopened natural-shift target remains an open empirical challenge. The retrospective Camelyon17 diagnostic in \KBSuppRef{app:camelyon17}{Appendix}{B.5} similarly exhibits conservative retention (zero \adapt decisions under an outcome-blind decision rule).
  \item \textbf{Retrospective Benchmarks vs.\ Prospective Deployment:} The controlled CIFAR-10-C and ImageNet-C
  experiments evaluate retrospective, opened panels under their track-specific declared protocols (including Protocol A and Protocol B). While these report empirical
  commitment and inclusion diagnostics on opened records, prospective evaluation on unopened environments is
  required to assess distribution-transfer properties.
  \item \textbf{Checkpoint and architecture transfer.} The primary CIFAR-10-C panel uses one
  source checkpoint. Additional checkpoint comparisons are descriptive: B6's clipped calibration
  rank does not establish nominal coverage, and incomplete execution provenance limits independent
  replication claims. Prospective transfer to unseen model families and environments remains
  untested (\KBSuppRef{app:model-replication}{Appendix}{the checkpoint comparison and its provenance qualification}).
  \item \textbf{External method comparisons.} The shared-feature baselines and scoped native
  components isolate specific decision rules, but do not constitute complete reproductions of
  all published learners or a full official-method comparison
  (\KBSuppRef{app:baseline-audit}{Appendix}{the supplementary baseline audit}).

\end{enumerate}
The controlled studies
quantify candidate-specific utility and commitment tradeoffs. CCT-20 provides retention evidence without
ADAPT exposure. These experiments do not establish population coverage or protection on unseen environments.
A unified statistical synthesis matrix consolidating sample sizes, evaluation units, nominal confidence intervals, multiplicity adjustments, and verdicts across the principal empirical tracks summarized here is reported in \KBSuppRef{tab:master-statistical-synthesis}{Table}{the master statistical synthesis table}.

\subsection{Operational Scope and Potential Harm}

A rejected update can prevent adaptation harm, but retaining a frozen model is not proof that the
model is accurate or suitable for a high-stakes setting. Misinterpreting an abstention as a
certificate would be particularly risky. In safety-critical deployments (e.g., automated diagnostics or autonomous navigation), a FREEZE or ABSTAIN action should not be assumed to guarantee safety; rather, it should function as an automated fallback trigger, preserving the frozen baseline while routing ambiguous cases to human clinical review or secondary fail-safe systems.
The reported experiments do not validate clinical, safety-critical, or autonomous operation.
No held-out physical camera deployment or common end-to-end latency result is claimed.

An unavailable assessment must remain distinct from evidence of negative benefit: both retain
the frozen model, but only the latter supports a strict FREEZE decision. The local integration
checks this behavior; its tests do not establish calibration validity or field reliability.
Detailed implementation and reproducibility information is provided in the supplement.

\section{Conclusion}
\label{sec:compact-conclusion}

The update decision depends on what the available evidence can distinguish. For the declared
binary target class, K-Bound identifies the feasible population-benefit interval and its exact
strict-decision frontier. The evidence-fibre audit establishes the accompanying limit: unchanged
label-free observations cannot justify an arbitrarily tighter residual bound. The paired-transport
construction illustrates a separate multiclass route under an explicit conditional-shift budget
and sampling assumptions; its usefulness depends on the credibility of those assumptions.

KGA makes the empirical question concrete by predicting paired evaluation-cell benefit and
calibrating an interval around that prediction. Its value is candidate- and operating-point-dependent.
Controlled routing gains coexist with simpler-rule wins. In the shared-predictor comparison,
scaled calibration improves CIFAR Tent's harm-weighted loss, but the advantage narrows at matched
exposure, and nominal 90\% constant-cell inclusion falls to $50.51$--$60.69\%$ on held-out CIFAR
corruptions. Useful decisions and reliable intervals must therefore be evaluated separately.

Natural-shift evidence is more limited: CCT-20 demonstrates retention without ADAPT exposure,
and So2Sat stops before target scoring. The present studies do not establish population-risk
protection or useful selective routing on unseen natural environments. The next empirical test
is whether independently supported calibration can accept beneficial updates and reject harmful
ones under a fixed protocol. Where the available information cannot support either strict sign,
abstention remains the justified decision.
